\documentclass[11pt]{article}

\usepackage[margin=1in]{geometry}
\usepackage{amsmath,amssymb,amsthm,mathtools}
\usepackage{booktabs}
\usepackage{enumitem}
\usepackage[numbers,sort&compress]{natbib}
\usepackage{hyperref}
\usepackage{microtype}
\usepackage{xcolor}

\hypersetup{colorlinks=true,linkcolor=blue,citecolor=blue,urlcolor=blue}
\allowdisplaybreaks

\newtheorem{proposition}{Proposition}
\newtheorem{lemma}{Lemma}
\newtheorem{definition}{Definition}
\numberwithin{equation}{section}

\newcommand{\R}{\mathbb R}
\newcommand{\E}{\mathbb E}
\newcommand{\N}{\mathcal N}
\newcommand{\dd}{\,\mathrm d}
\newcommand{\tr}{\operatorname{tr}}
\newcommand{\diag}{\operatorname{diag}}
\newcommand{\vecop}{\operatorname{vec}}
\newcommand{\sgq}{\operatorname{SGQ}}
\newcommand{\chol}{\operatorname{chol}}
\newcommand{\ellhat}{\widehat\ell}
\newcommand{\bfone}{\mathbf 1}

\title{An Expanded Implementation-Ready Mathematical Companion To Fixed Sparse-Grid Quadrature Filtering And Its Analytical Gradient}
\author{Chak Wong}
\date{2026-06-06}

\begin{document}
\maketitle

\begin{abstract}
This note develops fixed sparse-grid quadrature filtering as an
implementation-ready high-dimensional filtering proposal for the regime where
exact nonlinear filtering is unavailable but a deterministic Gaussian-surrogate
value-and-gradient target is still scientifically useful.  Starting from the
sparse-grid quadrature nonlinear filter of Jia, Xin, and Cheng
\cite{jia2012sparsegrid}, the note reconstructs the Gaussian approximation and
sparse-grid formulas in source order, then defines a fixed approximate
likelihood scalar together with an analytical derivative of that same scalar on
a declared branch.  The document is written for both review and implementation:
it states the exact filtering target, the narrower carried Gaussian surrogate,
the fixed cloud-construction rule, the per-step value recursion, the
same-scalar branch contract, the analytical gradient recursion, and the
finite-difference validity rule.  The central limitation is explicit.  The
reported scalar is not the exact nonlinear filtering likelihood, and the method
does not preserve general non-Gaussian posterior shape; it is a deterministic
Gaussian-surrogate lane whose usefulness depends on whether one carried Gaussian
and a fixed sparse-grid moment engine are adequate for the intended regime.
\end{abstract}

\section{Reader Orientation, Scope, And How To Use This Note}
\label{sec:p33-reader-orientation}

This note is written for three readers at once:
\begin{enumerate}[leftmargin=0.28in]
  \item implementation engineers who need an auditable mathematical contract,
  \item mathematically trained first-year-PhD-level workers who must implement
  the algorithm correctly from the note, and
  \item general scientific reviewers who need to judge whether the proposal is
  coherent, bounded, and worthy of approval.
\end{enumerate}

The document therefore serves two roles at once.  It is a mathematical
companion note for one deterministic approximate filtering lane, and it is an
implementation handoff document for that lane.  It is not a general survey of
nonlinear filtering and it is not a full software engineering manual.

\paragraph{Background assumed.}
The reader is assumed to know multivariate Gaussian notation, basic
state-space-model notation, Jacobian-style derivative notation, and the use of a
Cholesky factor for a positive-definite covariance matrix.  The note does not
assume prior familiarity with fixed-branch derivative contracts, sparse-grid
cloud merging, or same-scalar finite-difference validation; those are defined
here.

\paragraph{What this note provides internally.}
The note defines:
\begin{enumerate}[leftmargin=0.28in]
  \item the exact filtering target and the narrower approximate object carried by
  FixedSGQF,
  \item the fixed sparse-grid cloud construction in standardized coordinates,
  \item the one-step value recursion and the accumulated deterministic
  approximate likelihood scalar,
  \item the fixed-branch analytical derivative of that same scalar,
  \item the same-scalar finite-difference validity rule, and
  \item worked numerical checks that an implementation can use as direct
  oracles.
\end{enumerate}

\paragraph{How to read the note.}
Sections~\ref{sec:p32-problem-lane}--\ref{sec:p32-one-page} state the problem,
the selected approximation lane, and the exact-versus-approximate object map.
Sections~\ref{sec:p31-sgq-reconstruction}, \ref{sec:p31-toy-grid}, and
\ref{sec:p31-value-path} define the fixed cloud and the value recursion.
Sections~\ref{sec:p31-same-scalar} and \ref{sec:p31-gradient} define the saved
branch and the analytical gradient.  Section~\ref{sec:p31-fd} states the
finite-difference audit.  Implementation workers should treat the equations,
contracts, boxed algorithms, and worked example as the authoritative coding
material.

\paragraph{What this note does not claim.}
The note does not claim exact nonlinear filtering, exact nonlinear likelihood
evaluation, or faithful preservation of arbitrary posterior multimodality,
skewness, or hard all-coordinate interaction structure.  It defines one
deterministic Gaussian-surrogate lane and makes its approximation boundary
explicit.

\tableofcontents

\section{The High-Dimensional Filtering Problem And The FixedSGQF Lane}
\label{sec:p32-problem-lane}

The scientific problem is high-dimensional nonlinear filtering.  In a general
state-space model, the exact Bayesian filter propagates a full probability law
through time.  Once the transition or observation map is nonlinear, that law is
usually no longer describable by a finite list of moments, and in high
dimension it becomes unrealistic to treat exact function-valued propagation as a
practical computational object.  A review panel should therefore not read this
report as an attempt to rescue exactness by notation.  The question is narrower
and more honest: what deterministic approximation lane is mathematically
coherent, transparent about its compromises, and still useful enough to deserve
serious consideration?

The lane proposed here is FixedSGQF, short for fixed sparse-grid quadrature
filtering.  It is motivated by a common tension in scientific filtering work.
On one side, full non-Gaussian methods can represent richer posterior geometry,
but they often carry heavier approximation objects and more complicated
smoothness or implementation burdens.  On the other side, one may still want a
deterministic approximate likelihood surrogate whose value is reproducible, whose
dependence on model parameters can be differentiated analytically on a declared
branch, and whose numerical behavior is inspectable without Monte Carlo noise.
FixedSGQF belongs to the second lane.

The proposal is not to store the full filtering density.  It is to store one
Gaussian surrogate at each time, to evaluate the nonlinear moments required by
that Gaussian projection with a deterministic sparse-grid rule, and then to
freeze the grid, merge, and covariance-factor branches so that the reported
gradient differentiates the same scalar that the reported value routine
evaluates.  The gain from that choice is transparency: every approximation stage
can be named, inspected, and tested directly.  The cost is equally explicit: a
single Gaussian carried object cannot preserve general multimodal or strongly
skewed posterior structure, even when the quadrature itself is accurate for the
moments it is asked to approximate.

This is why FixedSGQF should not be read as a weaker version of the Zhao--Cui
squared tensor-train construction.  The two proposals serve different goals.
FixedSGQF is the deterministic Gaussian-surrogate lane with an analytical
same-scalar gradient.  Zhao--Cui squared TT is the richer non-Gaussian
density-approximation lane.  A serious high-dimensional filtering program may
want both.  The purpose of this report is to show that the FixedSGQF lane is
self-contained, mathematically explicit, and implementable enough to stand on
its own merits.

For the review chair, the decision test is therefore the following.  This lane
is worth advancing if the program needs a deterministic, inspectable, and
analytically differentiable likelihood surrogate whose approximation boundary is
fully exposed, even while richer non-Gaussian density lanes remain necessary for
problems where one carried Gaussian is too weak.  The report must make that
trade-off explicit enough that approval does not depend on optimism or tacit
background knowledge.

\section{FixedSGQF In One Page}
\label{sec:p32-one-page}

FixedSGQF is built from three deliberate approximation decisions.

\paragraph{First approximation: Gaussian carried state.}
At time \(t\), the exact filter carries the full conditional law
\(p_\theta(x_t\mid y_{1:t})\).  FixedSGQF instead carries only one Gaussian,
recorded by a mean and covariance pair \((m_t,P_t)\).  This is the main
structural narrowing of the proposal.  The report must therefore keep separate
throughout what belongs to the exact Bayesian filter and what belongs only to
its Gaussian surrogate.

\paragraph{Second approximation: sparse-grid moment evaluation.}
Once the carried object has been narrowed to a Gaussian, the filter still needs
nonlinear prediction and observation moments.  FixedSGQF evaluates those moments
with a deterministic sparse-grid quadrature rule built in standard Gaussian
coordinates and then mapped into physical state coordinates by the current
covariance factor.  The rule is fixed before likelihood evaluation, so the
moments are deterministic functions of the declared branch and the model
parameters.

\paragraph{Third approximation: fixed-branch differentiation.}
The scalar differentiated in this note is not the derivative of an adaptive grid
selection algorithm.  It is the derivative of a declared deterministic scalar
computed on a frozen branch.  Grid level, merged cloud, duplicate-node policy,
covariance-factor family, and veto rules are fixed first.  Within that branch,
the numerical values of the moments, gains, and likelihood contributions still
depend on \(\theta\), but the structural route through which they are computed
does not change.

The resulting scalar is
\begin{equation}
  \ellhat_T^{\rm FSGQ}(\theta)
  =
  \sum_{t=1}^T \ellhat_t^{\rm FSGQ}(\theta),
  \label{eq:p32-one-page-scalar}
\end{equation}
where each increment is the Gaussian innovation log likelihood implied by the
sparse-grid moment approximation.  The exact object and the approximate object
are therefore different in kind:
\begin{equation}
\begin{aligned}
  \text{exact object:}
  &\quad p_\theta(x_t\mid y_{1:t}), \\
  \text{approximate carried object:}
  &\quad \N(m_t,P_t), \\
  \text{accumulated deterministic scalar:}
  &\quad \ellhat_T^{\rm FSGQ}(\theta).
\end{aligned}
  \label{eq:p32-exact-vs-approximate-summary}
\end{equation}
That distinction is not bookkeeping.  It is the central scientific claim of the
report: FixedSGQF is worthwhile only if one judges that this narrower,
deterministic, Gaussian-projection lane is still scientifically useful in the
regime under study.

\paragraph{Approximation ladder.}
The rest of the report unfolds four commitments, each narrower than the one
before it:
\begin{enumerate}[leftmargin=0.28in]
  \item \textbf{Exact target:} the scientific target is the exact nonlinear
  filtering law \(p_\theta(x_t\mid y_{1:t})\).
  \item \textbf{Gaussian carried surrogate:} FixedSGQF carries only one
  Gaussian summary \(\N(m_t,P_t)\) instead of the full law.
  \item \textbf{Sparse-grid moment closure:} the moments defining that
  Gaussian surrogate are approximated by a deterministic sparse-grid rule.
  \item \textbf{Fixed-scalar derivative contract:} the reported gradient is the
  derivative of the declared deterministic surrogate scalar on a frozen branch,
  not the derivative of a live adaptive algorithm.
\end{enumerate}
A chair should be able to retain this ladder after one reading: exact target,
Gaussian surrogate, sparse-grid moment closure, fixed-scalar derivative
contract.  Every later section exists to unpack one rung of this scheme.

\paragraph{Operational summary for implementers.}
A first implementation should preserve the following runtime order.
\begin{enumerate}[leftmargin=0.28in]
  \item Build and save one fixed sparse-grid cloud in standardized Gaussian
  coordinates.
  \item Start from an initial Gaussian surrogate \((m_0,P_0)\).
  \item At each time step, place the previous cloud with the current factor,
  push points through the transition, and compute predictive moments.
  \item If the predictive covariance fails the declared branch, stop and report
  a predictive-branch failure.
  \item Place the predictive cloud, push points through the observation map,
  and compute \(\bar z_t\), \(S_t\), and \(C_{xz,t}\).
  \item If the innovation covariance fails the declared branch, stop and report
  an innovation-branch failure.
  \item Form the one-step Gaussian innovation log likelihood, update the running
  scalar, and update the carried Gaussian.
  \item If gradients are requested, propagate the same fixed-branch derivative
  recursion on the same cloud and under the same branch policies.
\end{enumerate}
This summary is not a replacement for the later derivation.  It is the short
execution picture a junior implementer should retain before reading the longer
sections.

\section{Approximation Hierarchy And Coordinate Walk}
\label{sec:p32-hierarchy}

The method in this note is easiest to read if those four commitments are kept in
the same order in which a program would apply them.  Start with the exact
Bayesian filter.  At time \(t\), the exact predictive density is a function
\[
  x_t\mapsto p_\theta(x_t\mid y_{1:t-1}),
\]
and the exact filtering density after seeing \(y_t\) is
\[
  x_t\mapsto p_\theta(x_t\mid y_{1:t}).
\]
Those are full functions on \(\R^{n_x}\).  A general nonlinear observation can
make them skewed, curved, or multimodal.  FixedSGQF does not store those
functions.  It replaces each carried density by one Gaussian:
\[
  p_\theta(x_t\mid y_{1:t})
  \quad\hbox{is represented only by}\quad
  (m_t,P_t).
\]
The first approximation is therefore a Gaussian projection.  The second
approximation is quadrature: the Gaussian moments needed for the projection are
computed by a deterministic sparse-grid sum.  The third approximation is the
fixed-branch contract: the sparse-grid cloud, duplicate merge, covariance
factor branch, and failure exits are held fixed so that the analytical
derivative differentiates the same scalar that the value routine evaluates.

The coordinates are:
\begin{align}
  \xi &\in \R^{n_x},
  &\xi&\sim\N(0,I_{n_x})
  &&\hbox{standard quadrature coordinate},
  \label{eq:p32-standard-coordinate}\\
  x &= m+C\xi,
  &CC^\top&=P
  &&\hbox{physical state coordinate under a carried Gaussian},
  \label{eq:p32-physical-coordinate}\\
  a &= f_\theta(x),
  &&
  &&\hbox{transition image used for prediction},
  \label{eq:p32-transition-coordinate}\\
  z &= h_\theta(x),
  &&
  &&\hbox{observation image used for the innovation}.
  \label{eq:p32-observation-coordinate}
\end{align}
The sparse grid is built once in \(\xi\)-space.  The same standardized nodes
\(\xi^{(r)}\) are then placed into the physical state space by the current
Gaussian factor \(C\).  This is why the same saved cloud can be reused at every
time while the physical points move as \(m_t\), \(P_t\), and \(\theta\) move.

The likelihood scalar is also an approximation hierarchy.  The exact
predictive likelihood factor is
\begin{equation}
  p_\theta(y_t\mid y_{1:t-1})
  =
  \int g_\theta(y_t\mid x_t)
  p_\theta(x_t\mid y_{1:t-1})\dd x_t.
  \label{eq:p32-exact-evidence-factor}
\end{equation}
FixedSGQF instead forms a Gaussian innovation model
\begin{equation}
  Y_t\mid y_{1:t-1}
  \approx
  \N(\bar z_t,S_t),
  \label{eq:p32-gaussian-innovation-model}
\end{equation}
where \(\bar z_t\) and \(S_t\) are sparse-grid approximations of the moments of
\(h_\theta(X_t^-)+\varepsilon_t\).  The scalar differentiated later is
\begin{equation}
  \ellhat_T^{\rm FSGQ}(\theta)
  =
  \sum_{t=1}^T
  \log \N(y_t;\bar z_t(\theta),S_t(\theta)).
  \label{eq:p32-hierarchy-scalar}
\end{equation}
Equation \eqref{eq:p32-hierarchy-scalar} is useful because it is deterministic,
reproducible, and differentiable on a fixed branch.  It is narrower than the
exact likelihood in \eqref{eq:p32-exact-evidence-factor}.  This distinction is
not a weakness in the derivation; it is the contract that makes the method
implementable and auditable.

\paragraph{Why moderate sparse-grid level can be plausible.}
Suppose the observation map is not an arbitrary fully coupled function of all
state coordinates, but has a low effective interaction order.  A typical
block-local observation has the form
\begin{equation}
  h(x)
  =
  \sum_{j=1}^{B} h_j(x_{\mathcal I_j})
  +\sum_{(j,k)\in\mathcal E} h_{jk}(x_{\mathcal I_j},x_{\mathcal I_k}),
  \label{eq:p32-low-interaction-map}
\end{equation}
where each block \(\mathcal I_j\) is small and the edge set \(\mathcal E\) is
sparse.  The Gaussian moments of \eqref{eq:p32-low-interaction-map} depend
strongly on low-order coordinate interactions and weakly on most high-order
interactions.  A sparse-grid rule of modest level is designed to spend points
on lower-order tensor interactions before it spends points on the full tensor
product.  That is the mathematical reason the point count can be plausible.
It is not a theorem that the posterior is easy.  If the likelihood has
important all-coordinate interactions, narrow ridges, or separated modes, a
fixed moderate sparse grid can fail even though its point count is attractive.

\paragraph{A physical analogy without changing the mathematics.}
In a chemical kinetics model, a measured spectrum or reaction rate often
depends strongly on a few combinations of concentrations and rate constants,
and only weakly on many nuisance directions.  A deterministic sparse-grid rule
can be plausible in such a setting because it samples low-order coordinate
interactions systematically.  The filter still checks the result by innovation
diagnostics, positive-definiteness, level ladders, and finite differences.  The
argument is "this interaction structure makes a moderate grid worth trying,"
not "chemistry makes the approximation exact."

\section{What This Note Computes}
\label{sec:p31-computes}

This section is the report's first exact scope statement.  It answers three
questions that must remain distinct throughout the rest of the document:
what object is carried from time step to time step, what scalar is accumulated,
and what information about the full posterior law is deliberately discarded.

The object computed here is deliberately narrow.  At time \(t\), before
observing \(y_t\), the filter carries one Gaussian approximation
\begin{equation}
  X_{t-1}\mid y_{1:t-1}
  \approx
  \N(m_{t-1},P_{t-1}).
  \label{eq:p31-carried-gaussian}
\end{equation}
It then uses a deterministic sparse-grid quadrature rule to approximate the
Gaussian moments needed by a Kalman-style projection update.  The output after
observing \(y_t\) is again only
\begin{equation}
  X_t\mid y_{1:t}
  \approx
  \N(m_t,P_t),
  \label{eq:p31-output-gaussian}
\end{equation}
together with the approximate Gaussian innovation log likelihood
\begin{equation}
  \ellhat_t
  =
  -\frac12
  \left\{
  \log\det S_t+v_t^\top S_t^{-1}v_t+n_y\log(2\pi)
  \right\}.
  \label{eq:p31-innovation-loglik-opening}
\end{equation}
The fixed-SGQF proposal is to use the sum
\begin{equation}
  \ellhat_T^{\rm FSGQ}(\theta)=\sum_{t=1}^T \ellhat_t^{\rm FSGQ}(\theta)
  \label{eq:p31-total-fixed-scalar-opening}
\end{equation}
as a deterministic approximate likelihood scalar and to differentiate that
same scalar analytically.

This note therefore computes an approximate likelihood and a Gaussian
projection recursion.  It does not compute the full nonlinear posterior law.
Sparse-grid exactness is exactness of a quadrature formula for selected
Gaussian-weighted polynomial integrands.  It is not a guarantee that the
posterior is Gaussian, unimodal, or accurately represented by \((m_t,P_t)\).

\paragraph{A scalar cell that keeps us honest.}
Let
\begin{equation}
  X\sim\N(0,P),\qquad Y=X^2+\varepsilon,\qquad
  \varepsilon\sim\N(0,R).
  \label{eq:p31-quadratic-cell}
\end{equation}
For \(y>0\) and small \(R\), the exact posterior density is proportional to
\begin{equation}
  \exp\!\left[-\frac{x^2}{2P}
  -\frac{(y-x^2)^2}{2R}\right],
  \label{eq:p31-quadratic-posterior}
\end{equation}
which can have two separated lobes near \(+\sqrt y\) and \(-\sqrt y\).  A
Gaussian projection filter can compute accurate moments of \(X^2\), but after
the update it still stores only one Gaussian.  This is not a bug in the
quadrature rule.  It is the modeling choice in \eqref{eq:p31-output-gaussian}.

The same limitation explains the role of fixed SGQF in the high-dimensional
program.  It is a transparent fixed approximate likelihood target with a
well-defined gradient.  It is not the non-Gaussian density approximation lane;
that role belongs to squared tensor-train filtering.  The next section therefore
pins down the exact state-space problem to which this deterministic
Gaussian-surrogate lane is being applied, so that the later sparse-grid,
gradient, and comparison arguments refer to one common filtering target.

\section{State-Space Model And Exact Filtering Recursion}
\label{sec:p31-ssm}

Use the additive-noise state-space model
\begin{align}
  X_t &= f_\theta(X_{t-1})+\eta_t,
  &
  \eta_t&\sim\N(0,Q_\theta),
  \label{eq:p31-state-model}\\
  Y_t &= h_\theta(X_t)+\varepsilon_t,
  &
  \varepsilon_t&\sim\N(0,R_\theta),
  \label{eq:p31-observation-model}
\end{align}
with \(X_t\in\R^{n_x}\), \(Y_t\in\R^{n_y}\), and parameter \(\theta\in\R^p\).
The exact Bayesian prediction and update are
\begin{align}
  p_\theta(x_t\mid y_{1:t-1})
  &=
  \int p_\theta(x_t\mid x_{t-1})
       p_\theta(x_{t-1}\mid y_{1:t-1})\dd x_{t-1},
  \label{eq:p31-exact-prediction}\\
  p_\theta(x_t\mid y_{1:t})
  &=
  \frac{
  g_\theta(y_t\mid x_t)
  p_\theta(x_t\mid y_{1:t-1})}
  {\int g_\theta(y_t\mid x_t)
  p_\theta(x_t\mid y_{1:t-1})\dd x_t}.
  \label{eq:p31-exact-update}
\end{align}
These are the conceptual filtering equations used by Gaussian approximation
filters in Jia, Xin, and Cheng \cite[Section 2]{jia2012sparsegrid}.  In
nonlinear models, the integrals in \eqref{eq:p31-exact-prediction} and
\eqref{eq:p31-exact-update} usually do not close in finite-dimensional form.
Fixed SGQF replaces those exact densities by a Gaussian moment projection.

\subsection{Report-Wide Object Map And Implementation Conventions}
\label{subsec:p32-object-map}

The report uses the following principal objects throughout the value and
gradient paths.  An implementation worker should treat this table as the core
object inventory for the note.
\begin{center}
\small
\begin{tabular}{p{0.20\linewidth}p{0.44\linewidth}p{0.22\linewidth}}
\toprule
symbol & meaning & shape \\
\midrule
\(\xi^{(r)}\) & standardized sparse-grid node in \(\R^{n_x}\) & \(n_x\) \\
\(w_r\) & accumulated signed weight after duplicate merge & scalar \\
\(m_t,P_t\) & carried Gaussian mean and covariance after update & \(n_x\), \(n_x\times n_x\) \\
\(m_t^-,P_t^-\) & predictive Gaussian mean and covariance before update & \(n_x\), \(n_x\times n_x\) \\
\(C_t, C_t^-\) & covariance factors with \(P_t=C_tC_t^\top\), \(P_t^-=C_t^-(C_t^-)^\top\) & \(n_x\times n_x\) \\
\(\chi_{t-1}^{(r)}\) & previous-state physical cloud point & \(n_x\) \\
\(a_t^{(r)}\) & transition image of \(\chi_{t-1}^{(r)}\) & \(n_x\) \\
\(\chi_t^{(r)}\) & predictive physical cloud point & \(n_x\) \\
\(z_t^{(r)}\) & observation image of \(\chi_t^{(r)}\) & \(n_y\) \\
\(\bar z_t\) & predicted observation mean & \(n_y\) \\
\(S_t\) & innovation covariance & \(n_y\times n_y\) \\
\(C_{xz,t}\) & state-observation cross-covariance & \(n_x\times n_y\) \\
\(v_t\) & innovation residual \(y_t-\bar z_t\) & \(n_y\) \\
\(K_t\) & Gaussian projection gain & \(n_x\times n_y\) \\
\(\ellhat_t^{\rm FSGQ}\) & one-step deterministic approximate log-likelihood increment & scalar \\
\bottomrule
\end{tabular}
\end{center}

Three implementation conventions are fixed report-wide.
\begin{enumerate}[leftmargin=0.28in]
  \item \textbf{Factor convention.}  Every covariance factor is the lower
  triangular Cholesky factor with positive diagonal, written
  \(C(P)=\chol(P)\), whenever the declared branch is valid.
  \item \textbf{Solve convention.}  Formulas may use \(S_t^{-1}\) for compact
  mathematics, but the implementation contract is to compute all such actions by
  linear solves against the declared Cholesky factor of \(S_t\), not by forming
  an explicit inverse.
  \item \textbf{Noise convention.}  Additive process noise enters analytically
  through the covariance term \(Q_\theta\) in predictive moments; additive
  observation noise enters analytically through \(R_\theta\) in the innovation
  covariance.  This report does not use state augmentation for these additive
  noises.
\end{enumerate}

When floating-point accumulation produces a covariance that is not exactly
symmetric, the report assumes explicit symmetrization by
\((P+P^\top)/2\) before branch checking.  If the resulting matrix is not
positive definite on the declared Cholesky branch, the branch is vetoed rather
than silently repaired.  Any jitter, clipping, or eigenvalue flooring policy
would define a different scalar than the one studied here.

\section{Gaussian Projection From Moments}
\label{sec:p31-gaussian-projection}

Assume the predictive state is represented by
\begin{equation}
  X_t^-\sim\N(m_t^-,P_t^-).
  \label{eq:p31-predictive-gaussian}
\end{equation}
Let
\begin{equation}
  Z_t=h_\theta(X_t^-)
  \label{eq:p31-predicted-observation}
\end{equation}
be the noiseless predicted observation.  The Gaussian projection needs
\begin{align}
  \bar z_t &= \E[Z_t],
  \label{eq:p31-zbar}\\
  S_t &= R_\theta+\E[(Z_t-\bar z_t)(Z_t-\bar z_t)^\top],
  \label{eq:p31-S}\\
  C_{xz,t}&=\E[(X_t^- -m_t^-)(Z_t-\bar z_t)^\top].
  \label{eq:p31-Cxz}
\end{align}
Given these moments, define
\begin{align}
  v_t&=y_t-\bar z_t,
  &
  K_t&=C_{xz,t}S_t^{-1},
  \label{eq:p31-v-K}\\
  m_t&=m_t^-+K_tv_t,
  &
  P_t&=P_t^- - K_tS_tK_t^\top.
  \label{eq:p31-projection-update}
\end{align}

\begin{proposition}[Affine projection]
\label{prop:p31-affine-projection}
Among estimators of \(X_t^-\) of the form \(a+B Y_t\), the minimum mean-square
estimator under the moments in
\eqref{eq:p31-zbar}--\eqref{eq:p31-Cxz} uses \(B=C_{xz,t}S_t^{-1}\).  Storing
the resulting mean and covariance as a Gaussian gives
\eqref{eq:p31-projection-update}.
\end{proposition}

\begin{proof}
Center \(\widetilde X=X_t^- -m_t^-\) and
\(\widetilde Y=Y_t-\bar z_t\).  The squared-error objective for
\(m_t^-+B\widetilde Y\) is
\begin{equation}
  J(B)=\E\|\widetilde X-B\widetilde Y\|^2
  =
  \tr(P_t^-)-2\tr(BC_{xz,t}^\top)+\tr(BS_tB^\top).
  \label{eq:p31-affine-objective}
\end{equation}
Differentiating gives \(\nabla_BJ=-2C_{xz,t}+2BS_t\).  If \(S_t\) is
nonsingular, the stationary point is \(B=C_{xz,t}S_t^{-1}\).  Substitution gives
the residual covariance \(P_t^- - C_{xz,t}S_t^{-1}C_{xz,t}^\top\), which equals
\eqref{eq:p31-projection-update}.
\end{proof}

The approximation has two layers.  First, the moments
\eqref{eq:p31-zbar}--\eqref{eq:p31-Cxz} may be approximated by quadrature.
Second, the resulting posterior information is compressed to one Gaussian.

\subsection{How This Matches The Jia--Xin--Cheng Gaussian Approximation Block}
\label{subsec:p32-jia-gaussian-block}

Jia, Xin, and Cheng first write the Gaussian approximation filter before they
introduce sparse grids \cite[Section 2.2]{jia2012sparsegrid}.  In their
additive-noise setting, the prediction moment equations have the same
structure as
\begin{align}
  m_t^-&=\E[f_\theta(X_{t-1})],
  \label{eq:p32-source-pred-mean}\\
  P_t^-&=
  Q_\theta+
  \E[(f_\theta(X_{t-1})-m_t^-)(f_\theta(X_{t-1})-m_t^-)^\top],
  \label{eq:p32-source-pred-cov}
\end{align}
where \(X_{t-1}\sim\N(m_{t-1},P_{t-1})\) under the carried Gaussian
approximation.  The update moment equations have the structure
\begin{align}
  \bar z_t&=\E[h_\theta(X_t^-)],
  \label{eq:p32-source-zbar}\\
  S_t&=
  R_\theta+
  \E[(h_\theta(X_t^-)-\bar z_t)(h_\theta(X_t^-)-\bar z_t)^\top],
  \label{eq:p32-source-S}\\
  C_{xz,t}&=
  \E[(X_t^- -m_t^-)(h_\theta(X_t^-)-\bar z_t)^\top].
  \label{eq:p32-source-Cxz}
\end{align}
Then the Gaussian projection update is
\begin{align}
  K_t&=C_{xz,t}S_t^{-1},
  \label{eq:p32-source-K}\\
  m_t&=m_t^-+K_t(y_t-\bar z_t),
  \label{eq:p32-source-m-update}\\
  P_t&=P_t^- -K_tS_tK_t^\top.
  \label{eq:p32-source-P-update}
\end{align}
The sparse-grid part of the paper answers only how to approximate the
expectations in \eqref{eq:p32-source-pred-mean}--\eqref{eq:p32-source-Cxz}.
It does not change the fact that the filter stores
\((m_t,P_t)\).  In BayesFilter notation, the replacement is
\begin{equation}
  \E[F(\xi)]
  \quad\leadsto\quad
  A_{b,L}[F]
  \quad\leadsto\quad
  \sum_{r=1}^{M_{b,L}}w_rF(\xi^{(r)}).
  \label{eq:p32-expectation-to-cloud}
\end{equation}
For prediction,
\[
  F_{\rm pred}(\xi)=
  f_\theta(m_{t-1}+C_{t-1}\xi),
\]
and for observation,
\[
  F_{\rm obs}(\xi)=
  h_\theta(m_t^-+C_t^-\xi).
\]
This is the exact point where the source SGQF construction enters the filter.
Everything before \eqref{eq:p32-expectation-to-cloud} is Gaussian projection;
everything after it is fixed sparse-grid implementation of the needed
Gaussian expectations.  Before formalizing the quadrature machinery, it helps to
see two small motivating cases that make the computational problem concrete.

\section{Why Quadrature Appears: Two Motivating Toy Cases}
\label{sec:p34-motivating-cases}

The next two sections become easier to read if the reader first sees the two
questions they are trying to answer:
\begin{enumerate}[leftmargin=0.28in]
  \item what integral must be approximated once a nonlinear Gaussian moment is
  needed, and
  \item why naive tensor-product point placement becomes expensive as dimension
  grows.
\end{enumerate}

\subsection{Toy case A: one nonlinear Gaussian moment}
\label{subsec:p34-scalar-motivation}

Start with a scalar Gaussian input and a simple nonlinear observation map.
Let
\begin{equation}
  X\sim\N(m,P),
  \qquad
  h(X)=X^2.
  \label{eq:p34-scalar-gaussian-model}
\end{equation}
Suppose we want the Gaussian-projection observation mean
\begin{equation}
  \bar z = \E[h(X)] = \E[X^2].
  \label{eq:p34-scalar-target-mean}
\end{equation}
Write \(X=m+C\xi\) with \(\xi\sim\N(0,1)\) and \(C^2=P\).  Then
\begin{equation}
  \bar z
  =
  \int_{\R} (m+C\xi)^2\phi(\xi)\dd\xi.
  \label{eq:p34-standardized-integral}
\end{equation}
This is the first key point: once the state has been standardized, the moment
problem becomes a standard-normal weighted integral.  The quadrature rule is
therefore not an extra modeling choice layered on top of nowhere.  It is a
numerical approximation to the Gaussian expectation in
\eqref{eq:p34-standardized-integral}.

In this scalar example, a three-point symmetric rule already reproduces the
first few Gaussian moments exactly, which is why rules such as
\(\{-\sqrt3,0,\sqrt3\}\) appear naturally later.  The purpose of the next
section is to formalize this one-dimensional Gaussian rule.

\subsection{Toy case B: a two-dimensional Gaussian placement picture}
\label{subsec:p34-2d-motivation}

Now move from one coordinate to two.  Use a simple two-dimensional Gaussian
state,
\begin{equation}
  X\sim\N(m,P),
  \qquad
  X\in\R^2,
  \label{eq:p34-2d-gaussian}
\end{equation}
and imagine either a two-dimensional linear-Gaussian state-space step or a
slightly nonlinear observation moment built from that same Gaussian input.  The
point of this example is geometric rather than inferential: it shows how the
same standardized cloud is placed and why tensor products grow quickly.

Let \(X=m+C\xi\) with \(\xi\sim\N(0,I_2)\).  If each standardized coordinate
uses the three-point rule \(\{-\sqrt3,0,\sqrt3\}\), then the direct tensor
product uses all pairs
\begin{equation}
  (\xi_1,\xi_2)
  \in
  \{-\sqrt3,0,\sqrt3\}\times\{-\sqrt3,0,\sqrt3\},
  \label{eq:p34-2d-tensor-points}
\end{equation}
so already in two dimensions the rule has \(3^2=9\) points.  In three
dimensions the same construction has \(3^3=27\) points; in ten dimensions it
has \(3^{10}\) points.  This is the second key point: the tensor-product rule is
simple and transparent, but its point count grows too quickly in higher
dimension.

Sparse-grid construction tries to keep the same standardized-coordinate logic
while avoiding the full tensor explosion.  It does that by combining lower-level
tensor rules with signed coefficients, then merging duplicate standardized
nodes.  The next two sections formalize that construction.

\section{One-Dimensional Gaussian Quadrature And Tensor Products}
\label{sec:p31-ghq}

The scalar toy case above shows why a Gaussian expectation becomes a
standard-normal weighted integral after standardization.  This section now makes
that one-dimensional rule explicit, and then shows how the same rule is lifted to
multiple coordinates by tensor products.

\paragraph{How to read this section with the toy cases in mind.}
Toy case A should be kept in mind whenever a one-dimensional rule
\(I_\ell[\psi]\) appears: there the integrand is the scalar nonlinear moment in
\eqref{eq:p34-standardized-integral}.  Toy case B should be kept in mind whenever
a tensor product appears: there the integrand lives on two standardized
coordinates and the point set is the direct Cartesian product in
\eqref{eq:p34-2d-tensor-points}.  The formal notation below is therefore not new
mathematics disconnected from the examples; it is the abstract version of those
two simple cases.

\subsection{What problem is equation \texorpdfstring{\eqref{eq:p31-univariate-rule}}{9.3} solving?}

In Toy case A, we needed the scalar nonlinear Gaussian moment
\begin{equation}
  \bar z
  =
  \int_{\R}(m+C\xi)^2\phi(\xi)\dd\xi.
  \label{eq:p36-scalar-target-repeat}
\end{equation}
The problem is simple to state: the integrand is nonlinear, so we want a small
weighted sum that approximates this standard-normal integral without evaluating
it analytically every time.  A one-dimensional standard-normal quadrature rule at
level \(\ell\) is exactly such a weighted sum:
\begin{equation}
  I_\ell[\psi]
  =
  \sum_{a=1}^{s_\ell}\omega_{\ell,a}\psi(\zeta_{\ell,a})
  \approx
  \int_{\R}\psi(\xi)\phi(\xi)\dd \xi.
  \label{eq:p31-univariate-rule}
\end{equation}
Here:
\begin{itemize}[leftmargin=0.28in]
  \item \(\psi\) is the function we want to integrate,
  \item \(\zeta_{\ell,a}\) are the standardized evaluation points, or nodes,
  \item \(\omega_{\ell,a}\) are the corresponding weights, and
  \item \(s_\ell\) is the number of nodes used by level \(\ell\).
\end{itemize}
In Toy case A, the role of \(\psi\) is played by
\begin{equation}
  \psi(\xi)=(m+C\xi)^2.
  \label{eq:p35-scalar-psi}
\end{equation}
So equation \eqref{eq:p31-univariate-rule} just says: approximate the scalar
Gaussian moment in \eqref{eq:p36-scalar-target-repeat} by evaluating
\((m+C\xi)^2\) at a small list of standardized points and taking a weighted sum.

\paragraph{Why do nodes and weights exist at all?}
They are not arbitrary decoration.  They are chosen so that the weighted sum gets
important Gaussian moments right.  For example, a symmetric three-point rule can
be chosen so that it integrates constants, \(\xi^2\), and \(\xi^4\) exactly under
the standard-normal weight.  That is why the later moment-matching equations
matter: they explain where the concrete nodes and weights come from.

\subsection{A fully worked one-dimensional toy rule}

Take the symmetric three-point ansatz
\begin{equation}
  \{(-a,\omega),(0,\omega_0),(a,\omega)\}.
  \label{eq:p36-three-point-ansatz}
\end{equation}
To match the first relevant Gaussian moments, require
\begin{align}
  \omega_0+2\omega &= 1,
  \label{eq:p36-three-point-mass}\\
  2\omega a^2 &= 1,
  \label{eq:p36-three-point-second}\\
  2\omega a^4 &= 3.
  \label{eq:p36-three-point-fourth}
\end{align}
Dividing \eqref{eq:p36-three-point-fourth} by
\eqref{eq:p36-three-point-second} gives \(a^2=3\).  Then
\(\omega=1/6\) and \(\omega_0=2/3\).  So the rule becomes
\begin{equation}
  \left\{\left(-\sqrt3,\frac16\right),\left(0,\frac23\right),
  \left(\sqrt3,\frac16\right)\right\}.
  \label{eq:p36-three-point-rule}
\end{equation}
Returning to Toy case A, the scalar moment is then approximated by
\begin{equation}
  \bar z
  \approx
  \frac16(m-C\sqrt3)^2 + \frac23 m^2 + \frac16(m+C\sqrt3)^2.
  \label{eq:p36-scalar-three-point-approx}
\end{equation}
This is the simplest concrete instance of Gaussian quadrature used anywhere in
this note.

\subsection{What changes in two dimensions?}

Now move to Toy case B.  Instead of one standardized coordinate \(\xi\), we now
have two coordinates \((\xi_1,\xi_2)\).  Suppose we want to approximate a
Gaussian-weighted integral of a function \(F(\xi_1,\xi_2)\).  If we reuse the same
three-point rule in each coordinate, the direct tensor-product rule places all
pairs of one-dimensional nodes.  That gives the nine-point grid already listed in
\eqref{eq:p35-nine-point-grid}.

At this stage there is still no sparse-grid machinery.  We are only taking the
one-dimensional rule and using it in both coordinates.  The reader should think
of the tensor-product construction as a rectangular grid in standardized space.

\subsection{What is the multi-index \texorpdfstring{\(\mathbf i=(i_1,\ldots,i_b)\)}{i=(i1,...,ib)}?}

The multi-index arrives because different coordinates may use different
one-dimensional rule levels.  For example:
\begin{itemize}[leftmargin=0.28in]
  \item \((1,1)\) means ``use level 1 in coordinate 1 and level 1 in coordinate
  2,''
  \item \((1,2)\) means ``use level 1 in coordinate 1 and level 2 in coordinate
  2,''
  \item \((2,1)\) means ``use level 2 in coordinate 1 and level 1 in coordinate
  2,''
  \item \((2,2)\) means ``use level 2 in both coordinates.''
\end{itemize}
So the multi-index does not come from nowhere.  It is simply a bookkeeping device
that records which one-dimensional rule level is used in each coordinate.

Formally, for a multi-index \(\mathbf i=(i_1,\ldots,i_b)\), the corresponding
tensor rule is
\begin{equation}
  I_{\mathbf i}[F]
  =
  (I_{i_1}\otimes\cdots\otimes I_{i_b})[F].
  \label{eq:p31-tensor-rule}
\end{equation}
This means:
\begin{quote}
apply the one-dimensional rule of level \(i_1\) in the first coordinate, the
one-dimensional rule of level \(i_2\) in the second coordinate, and so on, then
take the full Cartesian product of the resulting nodes and weights.
\end{quote}

\subsection{A fully worked two-dimensional tensor example}

Take \(b=2\) and use:
\begin{equation}
  I_1:(0,1),
  \qquad
  I_2:\left\{\left(-\sqrt3,\frac16\right),\left(0,\frac23\right),
  \left(\sqrt3,\frac16\right)\right\}.
  \label{eq:p36-level1-level2-rules}
\end{equation}
Now work out three explicit tensor rules.

\paragraph{Rule \texorpdfstring{\((1,1)\)}{(1,1)}.}
Both coordinates use the one-point rule, so there is only one tensor point:
\begin{equation}
  (0,0)
  \quad\text{with weight}\quad 1\cdot 1 = 1.
  \label{eq:p36-rule-11}
\end{equation}

\paragraph{Rule \texorpdfstring{\((1,2)\)}{(1,2)}.}
The first coordinate uses level 1 and the second uses level 2.  So the tensor
points are:
\begin{equation}
  (0,-\sqrt3),\qquad (0,0),\qquad (0,\sqrt3),
  \label{eq:p36-rule-12-points}
\end{equation}
with weights
\begin{equation}
  \frac16,\qquad \frac23,\qquad \frac16.
  \label{eq:p36-rule-12-weights}
\end{equation}

\paragraph{Rule \texorpdfstring{\((2,1)\)}{(2,1)}.}
Now the first coordinate uses level 2 and the second uses level 1.  So the
tensor points are:
\begin{equation}
  (-\sqrt3,0),\qquad (0,0),\qquad (\sqrt3,0),
  \label{eq:p36-rule-21-points}
\end{equation}
with weights
\begin{equation}
  \frac16,\qquad \frac23,\qquad \frac16.
  \label{eq:p36-rule-21-weights}
\end{equation}

This is the concrete meaning of the tensor-rule notation.  The symbols
\(\xi^{\mathbf i,\alpha}\) and \(w^{\mathbf i,\alpha}\) introduced below are just
a compact way to name these explicit tensor points and their product weights.

Its nodes and weights are
\begin{align}
  \xi^{\mathbf i,\alpha}
  &=
  (\zeta_{i_1,\alpha_1},\ldots,\zeta_{i_b,\alpha_b}),
  \label{eq:p31-tensor-node}\\
  w^{\mathbf i,\alpha}
  &=
  \prod_{j=1}^b\omega_{i_j,\alpha_j},
  \qquad
  \alpha_j\in\{1,\ldots,s_{i_j}\}.
  \label{eq:p31-tensor-weight}
\end{align}
In Toy case B with two coordinates and the same three-point rule in each
coordinate, this abstract tensor product is exactly the nine-point Cartesian
product already listed in \eqref{eq:p34-2d-tensor-points}.  The notation
\(\xi^{\mathbf i,\alpha}\) is just the formal way to say ``take one selected
standardized point from coordinate 1 and one selected standardized point from
coordinate 2, then pair them.''

It helps to see the 2D case written out explicitly.  If both coordinates use the
three-point rule \(\{-\sqrt3,0,\sqrt3\}\) with weights
\(\{1/6,2/3,1/6\}\), then the tensor-product rule places the nine points
\begin{equation}
  (-\sqrt3,-\sqrt3),
  (-\sqrt3,0),
  (-\sqrt3,\sqrt3),
  (0,-\sqrt3),
  (0,0),
  (0,\sqrt3),
  (\sqrt3,-\sqrt3),
  (\sqrt3,0),
  (\sqrt3,\sqrt3),
  \label{eq:p35-nine-point-grid}
\end{equation}
with weights given by products of the one-dimensional weights.  For example,
\((0,0)\) has weight \((2/3)(2/3)=4/9\), while
\((\sqrt3,-\sqrt3)\) has weight \((1/6)(1/6)=1/36\).  Nothing mysterious has
happened yet: tensor products just place all combinations of one-dimensional
points.

A full tensor-product rule is simple but expensive: if every coordinate uses
\(s\) nodes, then \(M=s^b\).  The two-dimensional toy case above already shows
this growth geometrically.  With three points per coordinate, two dimensions use
nine points, three dimensions use twenty-seven points, and ten dimensions use
\(3^{10}\) points.  Jia, Xin, and Cheng introduce the sparse-grid
combination to reduce this point count for fixed accuracy level
\cite[Section 3]{jia2012sparsegrid}.

\subsection{A concrete three-dimensional preview of why sparse grids help}
\label{subsec:p37-3d-preview}

The panel asked for one concrete three-dimensional case showing what sparse-grid
combination is trying to buy.  Here is the cleanest version.

Use the same one-dimensional rules as before:
\begin{equation}
  I_1:(0,1),
  \qquad
  I_2:\left\{\left(-\sqrt3,\frac16\right),\left(0,\frac23\right),
  \left(\sqrt3,\frac16\right)\right\}.
  \label{eq:p37-level1-level2-rules}
\end{equation}
If all three coordinates use level 2, then the full tensor-product rule
\((2,2,2)\) places
\begin{equation}
  3\times 3\times 3 = 27
  \label{eq:p37-27-points}
\end{equation}
standardized points.  That is the direct three-dimensional analogue of the
nine-point 2D grid.

Now compare that with the sparse-grid band for \(b=3\) and \(L=2\).  The
allowed multi-indices satisfy
\begin{equation}
  2\le |\mathbf i| \le 4.
  \label{eq:p37-3d-band}
\end{equation}
Since every coordinate level is at least 1, the only admissible three-dimensional
multi-indices are
\begin{equation}
  (1,1,1),
  \qquad
  (1,1,2),
  \qquad
  (1,2,1),
  \qquad
  (2,1,1).
  \label{eq:p37-3d-active-indices}
\end{equation}
So the sparse-grid rule does not use the full level-2 tensor product
\((2,2,2)\).  Instead it combines one one-point tensor rule and three
three-point tensor rules.  The corresponding Smolyak coefficients are
\begin{equation}
  c_{(1,1,1)}=-2,
  \qquad
  c_{(1,1,2)}=1,
  \qquad
  c_{(1,2,1)}=1,
  \qquad
  c_{(2,1,1)}=1.
  \label{eq:p37-3d-coeffs}
\end{equation}
So even in this concrete case the sparse-grid rule is already a signed
combination of lower-level tensor rules rather than a simple subset of the
27-point tensor product.

The raw point counts before duplicate merging are therefore:
\begin{align}
  (1,1,1) &\leadsto 1 \text{ point},
  \label{eq:p37-3d-count-111}\\
  (1,1,2) &\leadsto 3 \text{ points},
  \label{eq:p37-3d-count-112}\\
  (1,2,1) &\leadsto 3 \text{ points},
  \label{eq:p37-3d-count-121}\\
  (2,1,1) &\leadsto 3 \text{ points},
  \label{eq:p37-3d-count-211}
\end{align}
for a total of
\begin{equation}
  1+3+3+3=10
  \label{eq:p37-3d-raw-count}
\end{equation}
raw contributions before duplicate merging.  The point \((0,0,0)\) appears in
all four tensor components, so after duplicate merging the final cloud has even
fewer distinct nodes:
\begin{equation}
  (0,0,0),
  \qquad
  (\pm\sqrt3,0,0),
  \qquad
  (0,\pm\sqrt3,0),
  \qquad
  (0,0,\pm\sqrt3),
  \label{eq:p37-3d-final-nodes}
\end{equation}
which is only
\begin{equation}
  1+2+2+2 = 7
  \label{eq:p37-3d-final-count}
\end{equation}
distinct points.

So in this concrete 3D case the sparse-grid construction reduces the direct
27-point tensor-product picture to a 7-point merged cloud.  Why does this work?
Because the low sparse-grid level keeps only low-order tensor interactions.
Here, level 2 in three dimensions keeps the center point and one-axis excursions,
but it does not yet spend points on all pairwise or all three-way corner
combinations.  This is exactly the intended economy of the sparse-grid rule: pay
for lower-order interaction structure before paying for the full tensor product.

The next section now formalizes this economy in general notation.

\section{Sparse-Grid Rule Reconstructed In Source Order}
\label{sec:p31-sgq-reconstruction}

The sparse-grid construction answers a specific question raised by the toy cases.
The scalar example shows why nonlinear Gaussian moments become standard-normal
weighted integrals after standardization.  The two-dimensional example shows why
naive tensor products grow quickly even when the placement logic is simple.  The
purpose of sparse-grid combination is to keep the standardized-coordinate moment
logic while reducing the point count relative to the full tensor product.

\paragraph{How to read this section with the toy cases in mind.}
Keep returning mentally to the worked low-dimensional examples from the previous
section.  In two dimensions, we explicitly listed the tensor rules
\((1,1)\), \((1,2)\), and \((2,1)\), and we saw the resulting points.  In three
dimensions, we now also have the concrete preview in which the full 27-point
tensor-product picture is reduced to a 7-point merged cloud.  This section asks
how to turn those concrete constructions into one scalable general rule.

\subsection{Start with the scalar and 2D cases before the general notation}

In one dimension, there is no sparse-grid combination to explain.  One simply
chooses a one-dimensional Gaussian rule and uses it.  The scalar toy problem is
therefore the basic building block.

In two dimensions, however, one already has several tensor rules available.  For
example, with the level-1 and level-2 rules from
\eqref{eq:p36-level1-level2-rules}, the tensor rules \((1,1)\), \((1,2)\), and
\((2,1)\) are all explicit point sets.  A sparse-grid rule will not keep all
possible tensor levels equally.  Instead, it chooses a structured subset of
those tensor rules and combines them with signed coefficients.

The general notation below is just the scalable version of that statement.

Suppose the filter needs a Gaussian expectation
\begin{equation}
  \mathcal I_b[F]
  =
  \int_{\R^b}F(\xi)\phi_b(\xi)\dd\xi,
  \qquad
  \phi_b(\xi)=(2\pi)^{-b/2}\exp(-\|\xi\|^2/2).
  \label{eq:p32-gaussian-integral-target}
\end{equation}
In the scalar toy case, this is just the one-dimensional standard-normal integral
in \eqref{eq:p34-standardized-integral}.  In the 2D toy case, it is a
Gaussian-weighted integral over \((\xi_1,\xi_2)\), approximated first by tensor
rules such as \((1,1)\), \((1,2)\), and \((2,1)\).

A full tensor rule approximates \(\mathcal I_b[F]\) by evaluating \(F\) on
every Cartesian product of one-dimensional nodes.  If each coordinate uses
\(s\) points, the point count is \(s^b\).  Jia, Xin, and Cheng replace that
full tensor rule by a signed linear combination of lower-level tensor rules
\cite[Section 3.1]{jia2012sparsegrid}.  The signs are not an implementation
accident.  They are inclusion-exclusion coefficients: low-level tensor rules
overlap, and overlapping contributions must be subtracted before higher-level
contributions are added.

\subsection{The Jia--Xin--Cheng Index Sets}
\label{subsec:p32-source-index-sets}

Jia, Xin, and Cheng group tensor-product levels by an auxiliary integer \(q\)
\cite[Eq. 27]{jia2012sparsegrid}.  In BayesFilter notation, for dimension
\(b\),
\begin{equation}
  N_q^b
  =
  \left\{\mathbf i=(i_1,\ldots,i_b)\in\mathbb N^b:
  \sum_{j=1}^b i_j=b+q\right\},
  \qquad q\ge 0,
  \label{eq:p32-Nqb}
\end{equation}
and \(N_q^b=\emptyset\) for \(q<0\).  A member
\(\mathbf i\in N_q^b\) says: use the one-dimensional rule \(I_{i_j}\) in
coordinate \(j\), then take their tensor product.  The larger \(q\) is, the
more coordinates can have levels above one.

This is the abstract version of the explicit 2D examples from the previous
section.  When \(b=2\):
\begin{itemize}[leftmargin=0.28in]
  \item \((1,1)\) lies in the slice where the total index is \(2\),
  \item \((1,2)\) and \((2,1)\) lie in the slice where the total index is
  \(3\),
  \item \((2,2)\) would lie in the slice where the total index is \(4\).
\end{itemize}
So the set \(N_q^b\) is simply a way to group multi-indices by their total
level.  It is not a new probabilistic object; it is a bookkeeping device for
organizing which tensor rules are being combined.

For a declared sparse-grid level \(L\), the source sums over
\[
  q=L-b,\ L-b+1,\ldots,L-1.
\]
Since \(N_q^b\) is empty when \(q<0\), this formula also covers the case
\(L<b\).  If \(\mathbf i\in N_q^b\), then
\begin{equation}
  |\mathbf i|=\sum_{j=1}^b i_j=b+q.
  \label{eq:p32-index-sum-q}
\end{equation}
Thus the same index set can be written without \(q\) as
\begin{equation}
  \mathcal A_{b,L}
  =
  \{\mathbf i\in\mathbb N^b: L\le |\mathbf i|\le L+b-1\},
  \qquad |\mathbf i|=\sum_{j=1}^b i_j.
  \label{eq:p31-level-band}
\end{equation}
This is exactly the band used by the implementable fixed-cloud construction in
this note.

To keep the 2D toy case in view, take \(b=2\) and \(L=2\).  Then the admissible
band becomes:
\begin{itemize}[leftmargin=0.28in]
  \item \((1,1)\), because \(|(1,1)|=2\),
  \item \((1,2)\), because \(|(1,2)|=3\),
  \item \((2,1)\), because \(|(2,1)|=3\),
  \item but not \((2,2)\), because \(|(2,2)|=4\) lies outside the band for this
  choice of \(L\).
\end{itemize}
So in the 2D case the active band already recovers the three tensor rules that
will later build the five-point merged toy cloud.  This is the concrete meaning
of the abstract band definition.

\subsection{The Smolyak Coefficients}
\label{subsec:p32-smolyak-coefficients}

In the source formula, the coefficient attached to the \(q\)-slice is
\[
  (-1)^{L-1-q}\binom{b-1}{L-1-q}.
\]
Using \(q=|\mathbf i|-b\) from \eqref{eq:p32-index-sum-q}, this becomes
\begin{equation}
  c_{\mathbf i}
  =
  (-1)^{L+b-1-|\mathbf i|}
  \binom{b-1}{L+b-1-|\mathbf i|}.
  \label{eq:p31-smolyak-coeff}
\end{equation}
Therefore the sparse-grid quadrature approximation is
\begin{equation}
  A_{b,L}[F]
  =
  \sum_{\mathbf i\in\mathcal A_{b,L}}
  c_{\mathbf i}I_{\mathbf i}[F].
  \label{eq:p31-smolyak-rule}
\end{equation}
Equation \eqref{eq:p31-smolyak-rule} is the BayesFilter form of
\citet[Eq. 26--29]{jia2012sparsegrid}.  The coefficient
\eqref{eq:p31-smolyak-coeff} can be negative.  That fact later explains why
the filter must check covariance matrices instead of assuming every weighted
second-moment sum is automatically positive semidefinite.

In the concrete 2D case with \(L=2\), the active tensor rules are
\((1,1)\), \((1,2)\), and \((2,1)\).  Their coefficients turn out to be
\begin{equation}
  c_{(1,1)}=-1,
  \qquad
  c_{(1,2)}=1,
  \qquad
  c_{(2,1)}=1.
  \label{eq:p36-2d-smolyak-coeffs}
\end{equation}
So the sparse rule in this case is not ``keep some of the nine tensor-product
points and drop the others.''  It is
\begin{equation}
  A_{2,2}[F] = -I_{(1,1)}[F] + I_{(1,2)}[F] + I_{(2,1)}[F].
  \label{eq:p36-2d-smolyak-rule}
\end{equation}
This is the place where the construction first looks strange to many readers,
and it is worth stating the point slowly: the sparse-grid rule is a signed
combination of tensor rules, not a simple pruning of one big tensor grid.
Because the component tensor rules overlap, some repeated point contributions
must later be added together and some net contributions become smaller than the
raw tensor pieces suggest.

\subsection{Univariate Moment Matching}
\label{subsec:p32-univariate-moment-matching}

The one-dimensional rules \(I_\ell\) must be chosen before
\eqref{eq:p31-smolyak-rule} becomes a concrete point set.  Jia, Xin, and Cheng
construct univariate rules by matching moments of a standard normal
\cite[Section 3.2]{jia2012sparsegrid}.  If \(p_a\) and \(\omega_a\) are the
one-dimensional nodes and weights, the matching equations are
\begin{equation}
  M_j
  =
  \int_{-\infty}^{\infty}x^j\phi(x)\dd x
  =
  \sum_{a=1}^{N_u}\omega_a p_a^j,
  \qquad j=0,1,\ldots,J.
  \label{eq:p32-univariate-moment-match}
\end{equation}
The moments on the left are explicit:
\begin{equation}
  M_j=
  \begin{cases}
  0, & j\ \hbox{odd},\\
  (j-1)(j-3)\cdots 3\cdot 1, & j\ \hbox{even}.
  \end{cases}
  \label{eq:p32-normal-moments}
\end{equation}
For a symmetric three-point rule
\[
  \{(-a,\omega),(0,\omega_0),(a,\omega)\},
\]
the first three nontrivial equations are
\begin{align}
  \omega_0+2\omega&=1,
  \label{eq:p32-three-point-mass}\\
  2\omega a^2&=1,
  \label{eq:p32-three-point-second}\\
  2\omega a^4&=3.
  \label{eq:p32-three-point-fourth}
\end{align}
Dividing \eqref{eq:p32-three-point-fourth} by
\eqref{eq:p32-three-point-second} gives \(a^2=3\).  Then
\(\omega=1/6\) and \(\omega_0=2/3\).  This derivation is why the toy cloud
below uses nodes \(\{-\sqrt3,0,\sqrt3\}\) with weights
\(\{1/6,2/3,1/6\}\).  The same moment-matching logic extends to the higher
level symmetric rules discussed by \citet[Section 3.2]{jia2012sparsegrid}.

This is also the exact place where Toy case A connects back to the formal rule.
The scalar nonlinear moment in \eqref{eq:p34-standardized-integral} needs a
one-dimensional Gaussian rule.  The three-point rule just derived is the first
concrete example of such a rule.  Toy case B then reuses the same one-dimensional
rule coordinatewise before tensor products or sparse-grid combinations are even
considered.

\subsection{From Formula To Point Cloud}
\label{subsec:p32-cloud-construction}

Equation \eqref{eq:p31-smolyak-rule} is still not an implementation object.
The implementation does not evaluate a symbolic formula directly; it must turn
that formula into an explicit merged list of standardized nodes and weights.
The 2D case is the right place to understand this slowly.

\paragraph{Step 1: list the contributing tensor rules.}
For \(b=2\) and \(L=2\), the active rules are \((1,1)\), \((1,2)\), and
\((2,1)\), with coefficients from \eqref{eq:p36-2d-smolyak-coeffs}.

\paragraph{Step 2: expand each tensor rule into explicit points.}
\begin{itemize}[leftmargin=0.28in]
  \item Rule \((1,1)\) contributes only \((0,0)\), with tensor weight \(1\),
  then sparse-grid coefficient \(-1\).
  \item Rule \((1,2)\) contributes
  \((0,-\sqrt3)\), \((0,0)\), and \((0,\sqrt3)\), with tensor weights
  \(1/6\), \(2/3\), and \(1/6\), then sparse-grid coefficient \(+1\).
  \item Rule \((2,1)\) contributes
  \((-\sqrt3,0)\), \((0,0)\), and \((\sqrt3,0)\), with tensor weights
  \(1/6\), \(2/3\), and \(1/6\), then sparse-grid coefficient \(+1\).
\end{itemize}

\paragraph{Step 3: merge repeated standardized nodes.}
The node \((0,0)\) appears in all three tensor components.  Its accumulated
weight is therefore
\begin{equation}
  -1 + \frac23 + \frac23 = \frac13.
  \label{eq:p36-origin-merge}
\end{equation}
The four axis nodes each appear only once, so they keep weight \(1/6\).
The final merged cloud is exactly the five-point cloud presented later in the toy
example section.

Now return to the general notation.  Each tensor rule \(I_{\mathbf i}\) has
points and weights
\begin{align}
  \xi^{\mathbf i,\alpha}
  &=
  (\zeta_{i_1,\alpha_1},\ldots,\zeta_{i_b,\alpha_b}),
  \label{eq:p32-smolyak-node}\\
  w^{\mathbf i,\alpha}
  &=
  \prod_{j=1}^b \omega_{i_j,\alpha_j},
  \qquad
  \alpha_j\in\{1,\ldots,s_{i_j}\}.
  \label{eq:p32-smolyak-tensor-weight}
\end{align}
The sparse-grid contribution of this tensor point is
\begin{equation}
  c_{\mathbf i}w^{\mathbf i,\alpha}.
  \label{eq:p32-smolyak-point-contribution}
\end{equation}
The same standardized point can appear in several tensor components.  Jia,
Xin, and Cheng's Algorithm 1 says to add a new point if it has not appeared
before, and otherwise update the old weight \cite[Algorithm 1]{jia2012sparsegrid}.
The mathematical object is therefore the accumulated dictionary
\begin{equation}
  D[\xi]
  =
  \sum_{\substack{\mathbf i\in\mathcal A_{b,L}\\
  \alpha:\,\xi^{\mathbf i,\alpha}=\xi}}
  c_{\mathbf i}w^{\mathbf i,\alpha}.
  \label{eq:p31-node-dictionary}
\end{equation}
After merging duplicate standardized nodes, enumerate the nonzero entries as
\begin{equation}
  \mathcal C_{b,L}
  =
  \{(\xi^{(r)},w_r)\}_{r=1}^{M_{b,L}}.
  \label{eq:p31-fixed-cloud}
\end{equation}
This report adopts the following deterministic cloud-construction convention.

\begin{enumerate}[leftmargin=0.28in]
  \item Enumerate multi-indices \(\mathbf i\in\mathcal A_{b,L}\) in
  lexicographic order.
  \item Within each tensor rule \(I_{\mathbf i}\), enumerate the local index
  tuples \(\alpha\) lexicographically.
  \item For each candidate standardized node
  \(\xi^{\mathbf i,\alpha}\), search the current dictionary for an existing
  node \(\xi\) satisfying
  \begin{equation}
    \|\xi^{\mathbf i,\alpha}-\xi\|_\infty
    \le
    \tau_{\rm merge}
    :=10^{-12}\max\{1,\|\xi\|_\infty\}.
    \label{eq:p32-merge-rule}
  \end{equation}
  If such a node exists, accumulate the signed weight
  \(c_{\mathbf i}w^{\mathbf i,\alpha}\) into its stored weight; otherwise add a
  new dictionary entry.
  \item Remove entries whose accumulated absolute weight is below the declared
  numerical-zero threshold.
  \item Sort the surviving merged nodes lexicographically and store the final
  cloud in that order.
\end{enumerate}

For the fixed BayesFilter scalar, \(\mathcal C_{b,L}\), the univariate
rules, the merge tolerance, the zero-weight rule, the sorting convention, and
the accumulated signed weights are part of the target.  If two implementations
use the same level \(L\) but merge nearly-equal floating-point nodes
differently, they have not necessarily evaluated the same scalar.

The most important conceptual point is now visible from the worked 2D case.
The cloud used later by the filter is not a symbolic Smolyak expression and not
a raw tensor rule.  It is the merged, sorted list of standardized nodes and
accumulated signed weights that remains after all tensor components have been
expanded and duplicates have been combined.  That is why the later filter stores
\(\mathcal C_{b,L}\) rather than a list of symbolic index sets.

\paragraph{Implementation-facing pseudocode summary.}
The runtime object produced by this section is a stored cloud builder routine of
the form \texttt{build\_fixed\_sparse\_grid\_cloud(}$(b,L,\{I_\ell\},\tau_{\rm merge},\tau_{\rm zero})$\texttt{)}
that returns the sorted merged cloud \(\mathcal C_{b,L}\), its weight-total
diagnostic, and any signed-weight diagnostics used later in validation.  The
mathematics above is the authoritative definition of that routine.

\subsection{Exactness, Point Count, And The UKF Relation}
\label{subsec:p32-source-theorems}

Theorem 3.1 of \citet{jia2012sparsegrid} is a quadrature exactness statement.
If each univariate \(I_\ell\) is exact for univariate polynomials through
degree \(2\ell-1\), then the sparse-grid rule is exact for multivariate
polynomials through the corresponding total degree stated in that theorem.
The theorem supports this inference:
\begin{equation}
  F\ \hbox{polynomial in the exactness class}
  \quad\Longrightarrow\quad
  A_{b,L}[F]=\mathcal I_b[F].
  \label{eq:p32-exactness-meaning}
\end{equation}
It does not support the stronger inference
\[
  \hbox{nonlinear posterior is exactly represented by one Gaussian}.
\]
FixedSGQF uses \eqref{eq:p32-exactness-meaning} to justify deterministic
moment approximation, not to claim exact Bayesian filtering.

The point-count proposition in \citet[Proposition 3.1]{jia2012sparsegrid}
also has a precise scope.  For fixed level \(L\), the number of sparse-grid
points grows polynomially in dimension, with highest order \(L-1\) under the
source conditions.  The core counting reason is visible from
\eqref{eq:p32-Nqb}: in the largest relevant slice, at most \(L-1\) coordinates
can have levels above one.  The number of ways to choose those coordinates is
of order \(\binom{b}{L-1}\), which is polynomial in \(b\) when \(L\) is fixed.
This explains the storage advantage over full tensor products.  It does not
choose the correct level for a hard model.

The nestedness proposition in \citet[Proposition 3.2]{jia2012sparsegrid}
explains when lower-level point sets are contained in higher-level point sets.
For implementation, its practical message is simple: nested univariate rules
can reduce the number of distinct points after merging.  Non-nested rules can
still be used, but the duplicate merge will be weaker and the point count can
be larger.

Finally, \citet[Theorem 3.2]{jia2012sparsegrid} shows that the unscented
transform point set appears as a level-2 sparse-grid member under the source's
choice of one-point and three-point univariate rules.  Here this theorem is
used as a conceptual bridge.  It tells a reader that SGQF is not a mysterious
new filtering principle; it is a systematic sparse-grid way to build
deterministic Gaussian moment points.  The selected BayesFilter object remains
the fixed cloud \eqref{eq:p31-fixed-cloud}.

\subsection{A Deterministic Level Ladder Before Failure}
\label{subsec:p32-level-ladder}

For implementation work, "choose a level" is not enough.  Use a declared
deterministic ladder:
\begin{equation}
  L\in\{2,3,4,\ldots,L_{\max}\},
  \label{eq:p32-level-ladder}
\end{equation}
with a point budget \(M_{\max}\), a covariance branch policy, and a validation
set of pilot parameter values \(\Theta_{\rm pilot}\).  For each level, build
\(\mathcal C_{b,L}\), record \(M_{b,L}\), and run the construction diagnostics
and pilot filtering diagnostics.  The first level that satisfies
\begin{align}
  M_{b,L}&\le M_{\max},
  \label{eq:p32-ladder-budget}\\
  \min_{t,\theta\in\Theta_{\rm pilot}}\lambda_{\min}(S_t(\theta))
  &\ge \epsilon_S,
  \label{eq:p32-ladder-S}\\
  \max_{\theta\in\Theta_{\rm pilot}}
  \left|
  \ellhat_T^{(L)}(\theta)-\ellhat_T^{(L-1)}(\theta)
  \right|
  &\le \epsilon_\ell
  \quad\hbox{for }L>2
  \label{eq:p32-ladder-likelihood}
\end{align}
is eligible for a fixed-branch likelihood.  If no level satisfies the ladder,
FixedSGQF returns a design failure rather than silently using a convenient
grid.  The numerical constants \(\epsilon_S\) and \(\epsilon_\ell\) are
diagnostic thresholds, not mathematical proof parameters.  A minimal default
table is given later in Section~\ref{sec:p32-defaults}.

\section{A Toy Fixed Grid With Duplicate Merging}
\label{sec:p31-toy-grid}

This section makes the sparse-grid cloud concrete.  Take \(b=2\) and \(L=2\).
Use
\begin{equation}
  I_1:\quad (0,1),
  \qquad
  I_2:\quad
  \left\{(-\sqrt3,1/6),(0,2/3),(\sqrt3,1/6)\right\},
  \label{eq:p31-toy-univariate}
\end{equation}
which is the familiar three-point standard-normal rule at level 2.  The band
\(\mathcal A_{2,2}\) contains
\begin{equation}
  (1,1),\qquad (1,2),\qquad (2,1),
  \label{eq:p31-toy-indices}
\end{equation}
with coefficients
\begin{equation}
  c_{(1,1)}=-1,\qquad c_{(1,2)}=1,\qquad c_{(2,1)}=1.
  \label{eq:p31-toy-coefficients}
\end{equation}
Before merging, the node \((0,0)\) appears three times.  Its accumulated weight
is
\begin{equation}
  -1+\frac23+\frac23=\frac13.
  \label{eq:p31-toy-origin-weight}
\end{equation}
The four axis nodes each receive weight \(1/6\).  The merged cloud is
\begin{center}
\begin{tabular}{ccc}
\toprule
node \(\xi\) & accumulated weight & source components \\
\midrule
\((0,0)\) & \(1/3\) & \((1,1),(1,2),(2,1)\) \\
\((-\sqrt3,0)\) & \(1/6\) & \((2,1)\) \\
\((\sqrt3,0)\) & \(1/6\) & \((2,1)\) \\
\((0,-\sqrt3)\) & \(1/6\) & \((1,2)\) \\
\((0,\sqrt3)\) & \(1/6\) & \((1,2)\) \\
\bottomrule
\end{tabular}
\end{center}
The weights sum to one:
\begin{equation}
  \frac13+4\cdot\frac16=1.
  \label{eq:p31-toy-weight-total}
\end{equation}
This toy cloud illustrates two implementation rules that are easy to get
wrong.  First, duplicate nodes must be accumulated before the cloud is used.
Second, signed sparse-grid coefficients are applied before accumulation; a
future level or a different univariate family can produce negative accumulated
weights, so covariance matrices require explicit diagnostics.

Algorithmically, this toy example should be read as a dictionary update.  Start
from an empty map from standardized nodes to accumulated weights.  Insert the
contribution of each tensor component in \eqref{eq:p31-toy-indices} with its
signed coefficient from \eqref{eq:p31-toy-coefficients}.  The origin node
\((0,0)\) receives three signed contributions and is stored only once with the
final accumulated weight from \eqref{eq:p31-toy-origin-weight}.  The resulting
five-node dictionary is exactly the object that is later reused by both the
value and gradient paths; no tensor-component provenance is needed after the
merge is complete.

\section{FixedSGQF Filtering Value Path}
\label{sec:p31-value-path}

Fix a standardized sparse-grid cloud
\begin{equation}
  \mathcal C=\{(\xi^{(r)},w_r)\}_{r=1}^M,
  \qquad
  \xi^{(r)}\in\R^{n_x},
  \quad
  w_r\in\R.
  \label{eq:p31-cloud-for-filter}
\end{equation}
The cloud is frozen before likelihood evaluation.  The value path has two
quadrature stages at each time.

\paragraph{What this section gives the implementer.}
This section is the authoritative one-step value recursion.  A value routine
constructed from this note should take the current carried Gaussian,
current observation, fixed cloud, and declared model objects as input, then
return either an accepted one-step increment with updated \((m_t,P_t)\) or a
branch failure record.

\subsection{Prediction Stage}

Let \(P_{t-1}=C_{t-1}C_{t-1}^\top\) on the declared square-root branch.  Place
the previous Gaussian cloud:
\begin{equation}
  \chi_{t-1}^{(r)}
  =
  m_{t-1}+C_{t-1}\xi^{(r)}.
  \label{eq:p31-pred-placed-points}
\end{equation}
Push points through the transition:
\begin{equation}
  a_t^{(r)}=f_\theta(\chi_{t-1}^{(r)}).
  \label{eq:p31-transition-points}
\end{equation}
The predicted Gaussian is
\begin{align}
  m_t^-&=\sum_{r=1}^Mw_ra_t^{(r)},
  \label{eq:p31-pred-mean}\\
  P_t^-&=Q_\theta+\sum_{r=1}^Mw_r
  (a_t^{(r)}-m_t^-)(a_t^{(r)}-m_t^-)^\top.
  \label{eq:p31-pred-cov}
\end{align}
After \eqref{eq:p31-pred-cov}, the implementation replaces \(P_t^-\) by
\((P_t^-+P_t^{-\top})/2\) and checks the declared positive-definite branch.
For the differentiable lane in this note, there is no adaptive jitter, floor,
clipping, pivot, or eigenvalue repair inside the scalar.  If the check fails,
the value path returns a veto rather than quietly changing the scalar.

\subsection{Observation Stage}

Factor \(P_t^-=C_t^-(C_t^-)^\top\) on the declared branch and place the
predictive cloud:
\begin{equation}
  \chi_t^{(r)}
  =
  m_t^-+C_t^-\xi^{(r)}.
  \label{eq:p31-obs-placed-points}
\end{equation}
Evaluate the observation map:
\begin{equation}
  z_t^{(r)}=h_\theta(\chi_t^{(r)}).
  \label{eq:p31-obs-values}
\end{equation}
Then compute
\begin{align}
  \bar z_t&=\sum_{r=1}^Mw_rz_t^{(r)},
  \label{eq:p31-obs-mean}\\
  \Delta_t^{(r)}&=z_t^{(r)}-\bar z_t,
  \label{eq:p31-obs-centered}\\
  S_t&=R_\theta+\sum_{r=1}^Mw_r
  \Delta_t^{(r)}\Delta_t^{(r)\top},
  \label{eq:p31-obs-cov}\\
  C_{xz,t}&=\sum_{r=1}^Mw_r
  (\chi_t^{(r)}-m_t^-)\Delta_t^{(r)\top}.
  \label{eq:p31-cross-cov}
\end{align}
The innovation and log likelihood are
\begin{align}
  v_t&=y_t-\bar z_t,
  \label{eq:p31-innovation}\\
  \ellhat_t^{\rm FSGQ}
  &=
  -\frac12
  \left\{\log\det S_t+v_t^\top S_t^{-1}v_t+n_y\log(2\pi)\right\}.
  \label{eq:p31-fsgq-loglik}
\end{align}
If \(S_t\) is not positive definite on the declared branch, the value path
returns a veto.  If it passes, the Gaussian projection update is
\begin{align}
  K_t&=C_{xz,t}S_t^{-1},
  \label{eq:p31-kalman-gain}\\
  m_t&=m_t^-+K_tv_t,
  \label{eq:p31-filter-mean}\\
  P_t&=P_t^- - K_tS_tK_t^\top.
  \label{eq:p31-filter-cov}
\end{align}

\paragraph{Value-path pseudocode summary.}
The runtime order defined by this section is:
\begin{enumerate}[leftmargin=0.28in]
  \item symmetrize and factor \(P_{t-1}\),
  \item place previous-state points and propagate through \(f_\theta\),
  \item form \((m_t^-,P_t^-)\) and veto if the predictive branch fails,
  \item factor \(P_t^-\), place predictive points, and propagate through
  \(h_\theta\),
  \item form \(\bar z_t\), \(S_t\), and \(C_{xz,t}\), and veto if the
  innovation branch fails,
  \item form \(v_t\), \(\ellhat_t^{\rm FSGQ}\), \(K_t\), \(m_t\), and \(P_t\),
  then veto if the carried covariance branch fails,
  \item otherwise return the one-step increment, updated carried Gaussian, and
  any diagnostics cached for derivative reuse.
\end{enumerate}
This pseudocode summary does not replace the equations above; it states their
required execution order.

\section{A Worked One-Step Example With A Numeric Oracle}
\label{sec:p32-worked-example}

This section gives one compact example that can be read two ways.  First, it is
small enough for a chair to teach back without reconstructing the notation from
scratch.  Second, it gives an implementation engineer a concrete one-step oracle
against which a first prototype can be checked.

Use the one-dimensional transition-observation model
\begin{align}
  X_0 &\sim \N(m_0,P_0),
  &m_0&=0.5,
  &P_0&=1,
  \label{eq:p32-example-initial}\\
  X_1 &= f_\beta(X_0)+\eta_1,
  &f_\beta(x)&=\beta x,
  &\eta_1&\sim \N(0,Q),
  \label{eq:p32-example-transition}\\
  Y_1 &= h_\theta(X_1)+\varepsilon_1,
  &h_\theta(x)&=x^2,
  &\varepsilon_1&\sim \N(0,R),
  \label{eq:p32-example-observation}
\end{align}
with
\begin{equation}
  Q=0.25,
  \qquad
  R=1,
  \qquad
  y_1=2.
  \label{eq:p32-example-QRy}
\end{equation}
Differentiate with respect to the single parameter component \(\beta\).  For
this chosen component, the transition depends on \(\beta\) but the observation
map does not, so
\begin{equation}
  \partial_\beta h_\theta(x)=0.
  \label{eq:p32-example-hbeta-zero}
\end{equation}
Evaluate all quantities at \(\beta=1\), and use the positive scalar Cholesky
branch \(C(P)=\sqrt P\) with no repair.  Because the example yields
\(P_1^-=5/4>0\) and \(S_1=43/8>0\), both the value and derivative lanes stay on
that same declared branch.  Use the three-point standard-normal rule
\begin{equation}
  \mathcal C_{1,2}
  =
  \left\{
  \left(-\sqrt3,\frac16\right),
  \left(0,\frac23\right),
  \left(\sqrt3,\frac16\right)
  \right\}.
  \label{eq:p32-example-cloud}
\end{equation}
The same three-node cloud is reused without change in both the value and
derivative calculations, so the example can be taught back as one declared branch with one declared cloud across both lanes.

\paragraph{Value path.}
Since \(C_0=\sqrt{P_0}=1\), the previous-state placed points are
\begin{equation}
  \chi_0^{(r)}=m_0+C_0\xi^{(r)}
  \in
  \{0.5-\sqrt3,\ 0.5,\ 0.5+\sqrt3\}.
  \label{eq:p32-example-prev-points}
\end{equation}
At \(\beta=1\), the transition images satisfy
\begin{equation}
  a_1^{(r)}=f_\beta(\chi_0^{(r)})=\beta\chi_0^{(r)}=\chi_0^{(r)}.
  \label{eq:p32-example-transition-images}
\end{equation}
The predicted Gaussian moments are therefore
\begin{equation}
  m_1^-=
  \sum_rw_ra_1^{(r)}=0.5,
  \qquad
  P_1^-=Q+\sum_rw_r(a_1^{(r)}-m_1^-)^2=1.25=\frac54.
  \label{eq:p32-example-pred-moments}
\end{equation}
The declared Cholesky factor is
\begin{equation}
  C_1^-=\chol(P_1^-)=\sqrt{1.25}=\frac{\sqrt5}{2}.
  \label{eq:p32-example-pred-factor}
\end{equation}
Hence the predictive placed points are
\begin{equation}
  \chi_1^{(r)}=m_1^-+C_1^-\xi^{(r)}
  \in
  \left\{0.5-\frac{\sqrt{15}}{2},\ 0.5,\ 0.5+\frac{\sqrt{15}}{2}\right\}.
  \label{eq:p32-example-pred-points}
\end{equation}
Passing those points through \(h_\theta(x)=x^2\) gives
\begin{equation}
  z_1^{(r)}
  \in
  \left\{2.0635083269,\ 0.25,\ 5.9364916731\right\}.
  \label{eq:p32-example-z-values}
\end{equation}
The centered predictive and observation values are
\begin{equation}
  d_1^{(r)}=a_1^{(r)}-m_1^-
  \in
  \{-\sqrt3,\ 0,\ \sqrt3\},
  \label{eq:p32-example-d}
\end{equation}
and
\begin{equation}
  \Delta_1^{(r)}=z_1^{(r)}-\bar z_1
  \in
  \{0.5635083269,\ -1.25,\ 4.4364916731\}.
  \label{eq:p32-example-Delta}
\end{equation}
The observation moments and Gaussian update objects are
\begin{align}
  \bar z_1&=\sum_rw_r z_1^{(r)}=1.5=\frac32,
  \label{eq:p32-example-zbar}\\
  S_1&=R+\sum_rw_r(z_1^{(r)}-\bar z_1)^2=5.375=\frac{43}{8},
  \label{eq:p32-example-S}\\
  C_{xz,1}&=\sum_rw_r(\chi_1^{(r)}-m_1^-)(z_1^{(r)}-\bar z_1)=1.25=\frac54,
  \label{eq:p32-example-Cxz}\\
  v_1&=y_1-\bar z_1=0.5=\frac12,
  \label{eq:p32-example-v}\\
  K_1&=C_{xz,1}S_1^{-1}=\frac{10}{43}\approx0.2325581395,
  \label{eq:p32-example-K}\\
  m_1&=m_1^-+K_1v_1=\frac{53}{86}\approx0.6162790698,
  \label{eq:p32-example-m}\\
  P_1&=P_1^- -K_1S_1K_1^\top=\frac{165}{172}\approx0.9593023256.
  \label{eq:p32-example-P}
\end{align}
The one-step deterministic approximate log likelihood is
\begin{equation}
  \ellhat_1^{\rm FSGQ}
  =
  -\frac12\left\{\log S_1+v_1^2S_1^{-1}+\log(2\pi)\right\}
  \approx -1.7830736342.
  \label{eq:p32-example-ell}
\end{equation}

\paragraph{Same-branch derivative with respect to \texorpdfstring{\(\beta\)}{beta}.}
The same cloud and the same positive scalar Cholesky branch are reused in the
derivative lane.  This example is chosen so that the parameter enters only
through the transition, not through the observation map, which is why
\(\partial_\beta h_\theta(x)=0\).  Because \(m_0\) and \(P_0\) are fixed in \(\beta\), the
previous-state placed points satisfy
\begin{equation}
  \dot\chi_0^{(r)}=0.
  \label{eq:p32-example-prev-dot}
\end{equation}
From \eqref{eq:p31-dot-transition},
\begin{equation}
  \dot a_1^{(r)}
  =
  D_xf_\beta(\chi_0^{(r)})\dot\chi_0^{(r)}+\partial_\beta f_\beta(\chi_0^{(r)})
  =
  \chi_0^{(r)}
  \in
  \{0.5-\sqrt3,\ 0.5,\ 0.5+\sqrt3\}.
  \label{eq:p32-example-adot}
\end{equation}
Hence
\begin{equation}
  \dot m_1^-=
  \sum_rw_r\dot a_1^{(r)}=m_0=0.5,
  \qquad
  \dot d_1^{(r)}=\dot a_1^{(r)}-\dot m_1^-
  \in
  \{-\sqrt3,\ 0,\ \sqrt3\},
  \label{eq:p32-example-mdot-dotd}
\end{equation}
and
\begin{equation}
  \dot P_1^-=
  \partial_\beta(\beta^2P_0+Q)=2.
  \label{eq:p32-example-Pminus-dot}
\end{equation}
Because \(C_1^-=\sqrt{P_1^-}\), the declared Cholesky derivative is
\begin{equation}
  \dot C_1^-=
  \frac{\dot P_1^-}{2C_1^-}
  =
  \frac{2}{\sqrt5}
  \approx 0.8944271910.
  \label{eq:p32-example-Cdot}
\end{equation}
Therefore the predictive point sensitivities are
\begin{equation}
  \dot\chi_1^{(r)}=\dot m_1^-+\dot C_1^-\xi^{(r)}
  \in
  \{-1.0491933385,\ 0.5,\ 2.0491933385\}.
  \label{eq:p32-example-chidot}
\end{equation}
Since \(h_\theta(x)=x^2\) and \(\partial_\beta h_\theta(x)=0\),
\begin{equation}
  \dot z_1^{(r)}=2\chi_1^{(r)}\dot\chi_1^{(r)}
  \in
  \{3.0143149884,\ 0.5,\ 9.9856850116\}.
  \label{eq:p32-example-zdot}
\end{equation}
The centered observation sensitivities are
\begin{equation}
  \dot\Delta_1^{(r)}=\dot z_1^{(r)}-\dot{\bar z}_1
  \in
  \{0.5143149884,\ -2,\ 7.4856850116\}.
  \label{eq:p32-example-Deltadot}
\end{equation}
The derived moment sensitivities are
\begin{align}
  \dot{\bar z}_1&=2.5=\frac52,
  &\dot v_1&=-2.5=-\frac52,
  \label{eq:p32-example-zbar-v-dot}\\
  \dot S_1&=14.5=\frac{29}{2},
  &\dot C_{xz,1}&=3.25=\frac{13}{4}.
  \label{eq:p32-example-S-Cxz-dot}
\end{align}
Let
\begin{equation}
  u_1=S_1^{-1}v_1=\frac{4}{43}\approx0.0930232558.
  \label{eq:p32-example-u}
\end{equation}
The three scalar ingredients in the score formula are
\begin{align}
  \tr(S_1^{-1}\dot S_1)&=\frac{\dot S_1}{S_1}=\frac{116}{43}\approx2.6976744186,
  \label{eq:p32-example-score-term1}\\
  2\dot v_1u_1&=-\frac{20}{43}\approx-0.4651162791,
  \label{eq:p32-example-score-term2}\\
  -u_1^2\dot S_1&=-\frac{232}{1849}\approx-0.1254732288.
  \label{eq:p32-example-score-term3}
\end{align}
Then the same-branch one-step score is
\begin{equation}
  \dot\ellhat_1^{\rm FSGQ}
  =
  -\frac12\left[
    \frac{\dot S_1}{S_1}
    +2\dot v_1u_1
    -u_1^2\dot S_1
  \right]
  =-\frac{1948}{1849}
  \approx -1.0535424554.
  \label{eq:p32-example-ell-dot}
\end{equation}
For later propagation,
\begin{equation}
  \dot K_1
  =
  \dot C_{xz,1}S_1^{-1}-K_1\dot S_1S_1^{-1}
  =-\frac{42}{1849}
  \approx -0.0227149811,
  \label{eq:p32-example-Kdot}
\end{equation}
so
\begin{align}
  \dot m_1&=\dot m_1^-+\dot K_1v_1+K_1\dot v_1
  =-\frac{343}{3698}
  \approx -0.0927528394,
  \label{eq:p32-example-mdot}\\
  \dot P_1&=\dot P_1^- -\dot K_1S_1K_1-K_1\dot S_1K_1-K_1S_1\dot K_1
  =\frac{2353}{1849}
  \approx 1.2725797729.
  \label{eq:p32-example-Pdot}
\end{align}
The example makes the report's main limitation concrete.  Every quantity above
is deterministic and internally checkable on the declared branch, yet the
observation law \(Y_1=X_1^2+\varepsilon_1\) still has the same qualitative
ability to generate non-Gaussian posterior structure emphasized earlier in the
scalar-cell warning.  FixedSGQF therefore gives a transparent Gaussian-surrogate
value-and-gradient oracle, not a full non-Gaussian posterior representation.

\paragraph{Implementation checklist for this oracle.}
A first implementation should be able to reproduce, up to declared floating-point
tolerance, at least the following quantities from this section:
\(m_1^-\), \(P_1^-\), \(\bar z_1\), \(S_1\), \(C_{xz,1}\), \(K_1\),
\(m_1\), \(P_1\), \(\ellhat_1^{\rm FSGQ}\), and
\(\dot\ellhat_1^{\rm FSGQ}\).  This is the compact end-to-end numerical oracle
for both the value path and the derivative path.

The next section now states the exact saved-scalar contract that the worked
example has instantiated numerically: which structural choices are frozen, which
numerical quantities are recomputed, and what it means for value, gradient, and
finite differences to refer to the same declared target.

\section{The Saved Scalar And Same-Scalar Contract}
\label{sec:p31-same-scalar}

The derivative in this note is meaningful only for the scalar actually
evaluated.  This section therefore does two things at once: it freezes the exact
structural choices that define the target, and it prepares the reader for why
the later derivative chapter must be read as the derivative of that declared
target rather than of a changing algorithm.  A fixed-SGQF branch is the tuple
\begin{equation}
\begin{aligned}
  \mathcal B
  =
  (&y_{1:T},\ \text{observation preprocessing},\
  m_0(\theta),P_0(\theta),\dot m_0,\dot P_0,\\
  &\mathcal C,\ \text{one-dimensional rules},\
  \text{duplicate-merge tolerance},\
  \text{zero-weight rule},\
  \text{node ordering},\\
  &C(P)=\chol(P),\
  \text{symmetrize-then-veto rule},\
  \epsilon_S,\epsilon_P,\\
  &f_\theta,h_\theta,Q_\theta,R_\theta,\
  D_xf_\theta,\partial_i f_\theta,
  D_xh_\theta,\partial_i h_\theta,\dot Q_\theta,\dot R_\theta).
\end{aligned}
  \label{eq:p31-branch-tuple}
\end{equation}
The fixed scalar is
\begin{equation}
  \ellhat_T^{\rm FSGQ}(\theta;\mathcal B)
  =
  \sum_{t=1}^T \ellhat_t^{\rm FSGQ}(\theta;\mathcal B),
  \label{eq:p31-fixed-scalar}
\end{equation}
where every term is computed by
\eqref{eq:p31-pred-placed-points}--\eqref{eq:p31-filter-cov}.  The same
\(\mathcal B\) must be used for values, gradients, and finite differences.
Changing the grid level, active index set, duplicate-merge tolerance,
zero-weight pruning rule, node ordering, Cholesky pivot, observation
preprocessing, initial-condition policy, eigenvalue floor, covariance clipping,
or point pruning changes the scalar.  The
derivative formulas below cover the open branch where the symmetrized
covariances pass the Cholesky and positive-definite checks without repair.  If
a jitter, floor, clipping, or pivoting map is desired, it must be declared as a
different scalar with its own derivative or stop rule.

\paragraph{Plain-language branch rule.}
The branch identity freezes structure, not numerical values.  When
\(\theta\) changes inside a same-scalar finite-difference check, the moments,
gains, and likelihood contributions are recomputed, but they must be recomputed
by the same cloud, the same merge/order policy, the same factor family, and the
same realized acceptance path through the declared branch logic.  This is the
operational meaning of ``same scalar'' for implementation review.

\paragraph{Frozen versus variable objects on a saved branch.}
The branch record \(\mathcal B\) freezes structural choices: the cloud,
merge/prune/order policy, factor family, veto thresholds, preprocessing
policy, and initial-condition policy.  It does \emph{not} freeze the numerical
values of the moments or filters.  Within a fixed branch,
\begin{equation}
\begin{aligned}
  \text{frozen structural objects:}
  &\quad \mathcal C,\ \tau_{\rm merge},\ \tau_{\rm zero},\ \text{ordering},
  \ C(\cdot),\ \epsilon_S,\epsilon_P,\\
  \text{parameter-dependent numerical objects:}
  &\quad m_t,P_t,C_t,m_t^-,P_t^-,C_t^-,\bar z_t,S_t,C_{xz,t},K_t,
  \ellhat_t^{\rm FSGQ}.
\end{aligned}
  \label{eq:p32-frozen-vs-variable}
\end{equation}
Thus same-scalar finite differences recompute the parameter-dependent numerical
objects by the same declared branch rule; they do not reuse old numerical
values.  In particular, \(m_0(\theta\pm he_i)\) and \(P_0(\theta\pm he_i)\)
are recomputed by the same initial-condition policy; only the policy is frozen,
not the numerical initial values.  What must remain unchanged is the structural
route by which those values are produced.

\paragraph{Operational branch identity for finite differences.}
For implementation review, define the branch-identity record
\begin{equation}
  \mathfrak B
  =
  (\mathcal C,\ \tau_{\rm merge},\ \tau_{\rm zero},\ \text{ordering},\
   \text{preprocessing policy},\ \text{initial-condition policy},\
   \text{symmetrize-then-veto rule},\ C(P)=\chol(P),\ \epsilon_S,\epsilon_P,
   \text{accepted/failure stage-time pattern}).
  \label{eq:p32-branch-identity-record}
\end{equation}
The accepted/failure stage-time pattern means the realized branch path through
each likelihood contribution: at every time step it records whether the
predictive covariance check, innovation covariance check, and carried-covariance
check were accepted, and if not, which stage failed first.  A central, forward,
or backward finite-difference row is branch-valid only if both perturbed
evaluations share the same saved structural metadata and the same realized
acceptance path as the base evaluation through every likelihood contribution
included in the comparison.  Equal factor family but different first failure
location counts as branch mismatch.  Equal first failure location but different
accepted-prefix pattern also counts as branch mismatch.  If either side changes
the accepted stage-time pattern, triggers a different veto location, or uses a
different cloud/merge/order/preprocessing/initial-condition specification, the
row is branch-invalid and is not interpreted as evidence about the derivative.

\section{Analytical Gradient Of The Fixed Scalar}
\label{sec:p31-gradient}

The value path has now fixed the declared scalar and the branch on which that
scalar is evaluated.  The next task is not to introduce a new target, but to
differentiate that same declared target in a way that can still be checked by
finite differences.  Differentiate with respect to parameter component \(\theta_i\).
A dot denotes \(\partial_i\), for example \(\dot m_t=\partial_i m_t\).  The data
\(y_t\) are recorded and do not depend on \(\theta\).

\paragraph{Why this derivative is operationally delicate.}
The one-step scalar at time \(t\) depends on objects constructed earlier in the
filter: the previous carried Gaussian, the covariance factor used to place the
cloud, the nonlinear transition and observation images, the innovation
covariance, and the posterior update.  The derivative must therefore propagate
through the entire recursion, and it is meaningful only when the value and
derivative paths stay on the same declared branch.

\paragraph{Current-score objects versus propagation objects.}
For the current likelihood contribution, the derivative closes once
\(\dot v_t\) and \(\dot S_t\) are known.  The derivative
\(\dot C_{xz,t}\) is not needed to close the current score, but it is needed to
form \(\dot K_t\) and hence to propagate \(\dot m_t\) and \(\dot P_t\) to the next
time step.  A correct implementation must keep these two roles separate.

\subsection{The Story Of The Derivative}
\label{subsec:p32-gradient-story}

The derivative is long because the scalar at time \(t\) depends on objects
created earlier in the filter.  This section therefore proceeds in six explicit
stages: factor derivative, point sensitivities, predictive moment
sensitivities, observation moment sensitivities, innovation score, and posterior
sensitivity propagation.  The chain is
\begin{equation}
\begin{aligned}
  \theta
  &\longmapsto
  (m_{t-1},P_{t-1})
  \longmapsto
  C_{t-1}
  \longmapsto
  \chi_{t-1}^{(r)}
  \longmapsto
  a_t^{(r)}
  \longmapsto
  (m_t^-,P_t^-)
  \longmapsto
  C_t^-                                      \\
  &\longmapsto
  \chi_t^{(r)}
  \longmapsto
  z_t^{(r)}
  \longmapsto
  (\bar z_t,S_t,C_{xz,t})
  \longmapsto
  (v_t,K_t,\ellhat_t,m_t,P_t).
\end{aligned}
  \label{eq:p32-gradient-chain}
\end{equation}
The sparse-grid nodes and weights do not appear with dots in this chain.  They
are saved branch objects.  Their job is to say where the deterministic moment
sum is evaluated, not to move when \(\theta\) moves.  The Gaussian factors
\(C_{t-1}\) and \(C_t^-\), however, do move with \(\theta\), because they are
functions of \(P_{t-1}\) and \(P_t^-\).  That is why the factor derivative is
part of the scalar contract.

For readers tracking the derivation operationally, the six stages below can be read as a
ledger of which object is introduced, what it depends on, and whether it closes the
present score or only prepares the next step.
\begin{equation}
\begin{aligned}
  	ext{factor derivative} &\Longrightarrow \dot C_{t-1},\dot C_t^- ,\
  	ext{prediction stage} &\Longrightarrow \dot\chi_{t-1}^{(r)},\dot a_t^{(r)},\dot m_t^-,\dot P_t^- ,\
  	ext{observation stage} &\Longrightarrow \dot\chi_t^{(r)},\dot z_t^{(r)},\dot{\bar z}_t,\dot v_t,\dot S_t,\dot C_{xz,t},\
  	ext{score stage} &\Longrightarrow \dot\ellhat_t^{
m FSGQ},\
  	ext{propagation stage} &\Longrightarrow \dot K_t,\dot m_t,\dot P_t .
\end{aligned}
  \label{eq:p32-derivation-ledger}
\end{equation}

The value path computes the scalar
\[
  \ellhat_t
  =
  -\frac12\{\log\det S_t+v_t^\top S_t^{-1}v_t+n_y\log(2\pi)\}.
\]
The gradient path differentiates the same line.  It is helpful to separate the
one-step derivative into the three pieces that will later reappear in the score
formula:
\begin{equation}
  \partial_i \ellhat_t
  =
  -\frac12\,
  \partial_i\log\det S_t
  -\frac12\,
  \partial_i(v_t^\top S_t^{-1}v_t)
  -\frac12\,\partial_i(n_y\log 2\pi).
  \label{eq:p32-score-three-pieces}
\end{equation}
Since the last term is zero, the current-period score is controlled entirely by
how \(S_t\) and \(v_t\) move.  That is why the gradient path needs only two
objects for the current likelihood contribution:
\[
  \dot v_t,\qquad \dot S_t.
\]
The cross-covariance derivative \(\dot C_{xz,t}\) is not needed for
\(\dot\ellhat_t\).  It is needed because \(C_{xz,t}\) forms the gain \(K_t\),
and \(K_t\) updates \(m_t\) and \(P_t\).  Those updated sensitivities become
the input to the next time step.  Thus the derivative has two roles:
\begin{align}
  \hbox{score role:}\quad
  &(\dot v_t,\dot S_t)\mapsto \dot\ellhat_t,
  \label{eq:p32-score-role}\\
  \hbox{propagation role:}\quad
  &(\dot C_{xz,t},\dot S_t,\dot v_t)
  \mapsto (\dot m_t,\dot P_t).
  \label{eq:p32-propagation-role}
\end{align}
Confusing these roles is a common implementation mistake.  A program can
appear to compute a correct one-step score while still propagating the wrong
next-step derivative.

\begin{center}
\small
\begin{tabular}{p{0.28\linewidth}p{0.18\linewidth}p{0.40\linewidth}}
\toprule
object & role & why it matters \\
\midrule
\(\dot v_t,\dot S_t\) & current score & sufficient for the present innovation score in \eqref{eq:p31-score} \\
\(\dot C_{xz,t}\) & propagation only & needed because \(K_t\) depends on \(C_{xz,t}\), even though the current score does not \\
\(\dot K_t\) & propagation only & transmits observation sensitivity into the posterior update \\
\(\dot m_t,\dot P_t\) & next-step state & become the differentiated inputs to the next prediction stage and the next factor derivative \\
\bottomrule
\end{tabular}
\end{center}

There are four places where the derivative can stop:
\begin{enumerate}[leftmargin=0.24in]
\item \(P_{t-1}\) or \(P_t^-\) is not on the saved factor branch.
\item \(S_t\) is not positive definite.
\item A finite-difference perturbation changes the saved branch, grid, merge,
or preprocessing policy.
\item A model derivative routine \(D_xf_\theta\), \(\partial_i f_\theta\),
\(D_xh_\theta\), \(\partial_i h_\theta\), \(\dot Q_\theta\), or
\(\dot R_\theta\) is missing or inconsistent with the value routine.
\end{enumerate}
On those events the note does not patch the scalar by hidden repair.  It
returns a branch failure.  This narrowness is exactly what lets the analytical
gradient be checked by finite differences.

\subsection{Square-Root Branch}

This first stage differentiates the factor map itself.  It is needed before any
point sensitivity can be written down, because the placed cloud points depend on
\(C_{t-1}\) and \(C_t^-\).  If this stage is wrong, every later derivative object
is wrong even if the score algebra itself is correct.  For the default
unpivoted Cholesky branch, \(P=LL^\top\) with positive diagonal and \(C(P)=L\).  The implementation convention is: always use the lower
triangular Cholesky factor returned by the declared factor map, never a
symmetric square root or pivoted alternative inside this scalar.  Given
symmetric \(\dot P\), set
\begin{equation}
  A=L^{-1}\dot P L^{-\top}.
  \label{eq:p31-chol-A}
\end{equation}
Define the lower-triangular operator
\begin{equation}
  \Phi(A)_{jk}
  =
  \begin{cases}
  A_{jk}, & j>k,\\
  A_{jj}/2, & j=k,\\
  0, & j<k.
  \end{cases}
  \label{eq:p31-chol-Phi}
\end{equation}
Then
\begin{equation}
  \dot L=L\Phi(A).
  \label{eq:p31-chol-derivative}
\end{equation}
Equation \eqref{eq:p31-chol-derivative} is the governing implementation recipe
for every factor derivative in this report.  In particular, whenever the value
path computes a covariance factor by \(C_t=\chol(P_t)\) or
\(C_t^-=\chol(P_t^-)\), the gradient path computes the matching
\(\dot C_t\) or \(\dot C_t^-\) by applying
\eqref{eq:p31-chol-A}--\eqref{eq:p31-chol-derivative} to the corresponding
\(\dot P_t\) or \(\dot P_t^-\).  No alternative factor derivative is allowed
inside the same scalar.
Indeed, \(\Phi(A)+\Phi(A)^\top=A\), so
\begin{equation}
  \dot LL^\top+L\dot L^\top
  =
  L\{\Phi(A)+\Phi(A)^\top\}L^\top
  =
  \dot P.
  \label{eq:p31-chol-check}
\end{equation}
If a different square-root map is used, the implementation must provide
\(\dot C=DC(P)[\dot P]\) for the same branch.  That alternative is outside the
scope of this report: every formula and algorithm below assumes the declared
lower-triangular Cholesky branch \(C(P)=\chol(P)\).

\subsection{Prediction Sensitivities}

This stage moves the derivative from the previous carried Gaussian into the
prediction cloud and the predictive moments.  Its outputs, \(\dot m_t^-\) and
\(\dot P_t^-\), are not yet enough for the current likelihood score, but they are
necessary because both the next observation cloud and the next factor derivative
depend on them.  From \eqref{eq:p31-pred-placed-points},
\begin{equation}
  \dot\chi_{t-1}^{(r)}
  =
  \dot m_{t-1}+\dot C_{t-1}\xi^{(r)}.
  \label{eq:p31-dot-pred-place}
\end{equation}
From \eqref{eq:p31-transition-points},
\begin{equation}
  \dot a_t^{(r)}
  =
  D_xf_\theta(\chi_{t-1}^{(r)})\dot\chi_{t-1}^{(r)}
  +\partial_i f_\theta(\chi_{t-1}^{(r)}).
  \label{eq:p31-dot-transition}
\end{equation}
The predicted mean derivative is
\begin{equation}
  \dot m_t^-=\sum_{r=1}^M w_r\dot a_t^{(r)}.
  \label{eq:p31-dot-pred-mean}
\end{equation}
Let
\begin{equation}
  d_t^{(r)}=a_t^{(r)}-m_t^-,
  \qquad
  \dot d_t^{(r)}=\dot a_t^{(r)}-\dot m_t^-.
  \label{eq:p31-pred-centered}
\end{equation}
Then
\begin{equation}
  \dot P_t^-
  =
  \dot Q_\theta
  +
  \sum_{r=1}^M w_r
  \left\{
  \dot d_t^{(r)}d_t^{(r)\top}
  +d_t^{(r)}\dot d_t^{(r)\top}
  \right\}.
  \label{eq:p31-dot-pred-cov}
\end{equation}

\subsection{Observation Sensitivities}

This stage converts predictive-state sensitivity into observation-side
sensitivity.  It is here that the derivative splits into the two objects needed
for the current score, \(\dot v_t\) and \(\dot S_t\), and the extra object
needed only for propagation, \(\dot C_{xz,t}\).  The logic is:
\begin{equation}
  \dot\chi_t^{(r)}
  \longmapsto
  \dot z_t^{(r)}
  \longmapsto
  (\dot{\bar z}_t,\dot v_t,\dot\Delta_t^{(r)})
  \longmapsto
  (\dot S_t,\dot C_{xz,t}).
  \label{eq:p32-observation-derivative-flow}
\end{equation}
Differentiate \eqref{eq:p31-obs-placed-points}:
\begin{equation}
  \dot\chi_t^{(r)}
  =
  \dot m_t^-+\dot C_t^-\xi^{(r)}.
  \label{eq:p31-dot-obs-place}
\end{equation}
Then
\begin{equation}
  \dot z_t^{(r)}
  =
  D_xh_\theta(\chi_t^{(r)})\dot\chi_t^{(r)}
  +\partial_i h_\theta(\chi_t^{(r)}).
  \label{eq:p31-dot-obs-value}
\end{equation}
The observation mean, residual, and centered residual derivatives are
\begin{align}
  \dot{\bar z}_t&=\sum_{r=1}^M w_r\dot z_t^{(r)},
  \label{eq:p31-dot-zbar}\\
  \dot v_t&=-\dot{\bar z}_t,
  \label{eq:p31-dot-v}\\
  \dot\Delta_t^{(r)}&=\dot z_t^{(r)}-\dot{\bar z}_t.
  \label{eq:p31-dot-Delta}
\end{align}
Differentiate \(S_t\) and \(C_{xz,t}\):
\begin{align}
  \dot S_t
  &=
  \dot R_\theta
  +
  \sum_{r=1}^M w_r
  \left\{
  \dot\Delta_t^{(r)}\Delta_t^{(r)\top}
  +\Delta_t^{(r)}\dot\Delta_t^{(r)\top}
  \right\},
  \label{eq:p31-dot-S}\\
  \dot C_{xz,t}
  &=
  \sum_{r=1}^M w_r
  \left\{
  (\dot\chi_t^{(r)}-\dot m_t^-)\Delta_t^{(r)\top}
  +(\chi_t^{(r)}-m_t^-)\dot\Delta_t^{(r)\top}
  \right\}.
  \label{eq:p31-dot-Cxz}
\end{align}

\subsection{Innovation Score}

At this point the current-period score problem is local.  Once \(\dot v_t\),
\(\dot S_t\), and the innovation solve variable are known, the pointwise cloud no
longer appears explicitly in the one-step score.  The proposition below is
therefore the score-stage payoff of all earlier derivative work.

\begin{proposition}[Fixed Gaussian innovation score]
\label{prop:p31-innovation-score}
Assume \(S_t\) is positive definite and the active branch is differentiable.
Let \(u_t\) solve
\begin{equation}
  S_tu_t=v_t.
  \label{eq:p31-score-solve}
\end{equation}
Then the derivative of \eqref{eq:p31-fsgq-loglik} is
\begin{equation}
  \dot\ellhat_t^{\rm FSGQ}
  =
  -\frac12
  \left[
  \tr(S_t^{-1}\dot S_t)
  +2\dot v_t^\top u_t
  -u_t^\top\dot S_tu_t
  \right].
  \label{eq:p31-score}
\end{equation}
\end{proposition}

\begin{proof}
Use
\begin{equation}
  \dd\log\det S=\tr(S^{-1}\dd S)
  \label{eq:p31-d-logdet}
\end{equation}
and
\begin{equation}
  \dd S^{-1}=-S^{-1}(\dd S)S^{-1}.
  \label{eq:p31-d-inverse}
\end{equation}
Therefore, with \(u=S^{-1}v\),
\begin{equation}
  \dd(v^\top S^{-1}v)
  =
  2(\dd v)^\top u-u^\top(\dd S)u.
  \label{eq:p31-d-quadratic}
\end{equation}
Substitute \(\dd S=\dot S_t\) and \(\dd v=\dot v_t\) into
\eqref{eq:p31-innovation-loglik-opening}.
\end{proof}

\subsection{Posterior Sensitivity Propagation}

If the note stopped at the innovation score, it would describe only how to
differentiate one isolated likelihood contribution.  The filter, however,
continues in time.  This stage therefore propagates the current observation-side
sensitivity back into the carried Gaussian that seeds the next step.  The
gradient must also propagate \(\dot m_t\) and \(\dot P_t\) into the next
time step.  The implementation order is:
\begin{equation}
  (\dot C_{xz,t},\dot S_t,\dot v_t)
  \longmapsto
  \dot K_t
  \longmapsto
  (\dot m_t,\dot P_t)
  \longmapsto
  \dot C_t
  \quad\hbox{at the next step through}\quad
  DC(P_t)[\dot P_t].
  \label{eq:p32-posterior-sensitivity-order}
\end{equation}
Differentiate \(K_t=C_{xz,t}S_t^{-1}\):
\begin{equation}
  \dot K_t
  =
  \dot C_{xz,t}S_t^{-1}
  -K_t\dot S_tS_t^{-1}.
  \label{eq:p31-dot-K}
\end{equation}
Then
\begin{align}
  \dot m_t
  &=
  \dot m_t^-+\dot K_tv_t+K_t\dot v_t,
  \label{eq:p31-dot-m}\\
  \dot P_t
  &=
  \dot P_t^-
  -\dot K_tS_tK_t^\top
  -K_t\dot S_tK_t^\top
  -K_tS_t\dot K_t^\top.
  \label{eq:p31-dot-P}
\end{align}
Equations \eqref{eq:p31-dot-pred-place}--\eqref{eq:p31-dot-P} define the
value-and-gradient recursion for the fixed scalar
\eqref{eq:p31-fixed-scalar}.  At this point the reader should retain one
structural summary: \(\dot v_t\) and \(\dot S_t\) close the current score, while
\(\dot C_{xz,t},\dot K_t,\dot m_t,\dot P_t\) exist because the filter must carry
its differentiated state forward in time.

\section{One Boxed Mathematical Algorithm}
\label{sec:p31-boxed-algorithm}

\fbox{\begin{minipage}{0.94\linewidth}
\textbf{FixedSGQF value and gradient for one fixed parameter component on one saved branch.}
The dot notation in this box means differentiation with respect to one declared
parameter component \(\theta_i\).  In particular,
\(\dot m_t,\dot P_t,\dot C_t^-,\dot S_t,\dot C_{xz,t}\), and
\(\dot\ellhat_t^{\rm FSGQ}\) in this box are the one-component versions of the
indexed quantities \(\dot m_t^{(i)},\dot P_t^{(i)},\dot C_t^{-(i)},\dot S_t^{(i)},
\dot C_{xz,t}^{(i)}\), and \(\partial_i\ellhat_t^{\rm FSGQ}\) used in the full
algorithm.  This boxed routine is a compact summary of the full
implementation-order algorithm in Section~\ref{sec:p32-end-to-end-algorithm};
when the two presentations differ in brevity, the longer end-to-end algorithm
is the authoritative execution order.

\begin{enumerate}[leftmargin=0.22in]
\item Choose dimension \(b=n_x\), level \(L\), one-dimensional rules
\(I_\ell\), merge tolerance, zero-weight threshold, and deterministic node
ordering.  Build and save the cloud \(\mathcal C=\{(\xi^{(r)},w_r)\}_{r=1}^M\)
by \eqref{eq:p31-level-band}--\eqref{eq:p31-fixed-cloud}, including duplicate
merge, near-zero pruning, and final deterministic sorting.  These cloud-construction
choices are part of the saved branch record and therefore part of the declared
scalar.
\item Save \(y_{1:T}\), observation preprocessing, \(m_0(\theta),P_0(\theta)\),
the one-component sensitivity policy for \(\theta_i\), the full model objects
\(f_\theta,h_\theta,Q_\theta,R_\theta\), the declared lower-triangular Cholesky
factor map \(C(P)=\chol(P)\), the symmetrize-then-veto rule, the eigenvalue thresholds
\eqref{eq:p32-failure-eigen}, and model derivative routines.  Together with the
cloud-construction choices in Step 1, these objects form the saved branch
record for the declared scalar.
\item Initialize \(m_0,P_0\), \(\dot m_0,\dot P_0\), the running scalar
\(\ellhat_T^{\rm FSGQ}\gets 0\), and the running derivative
\(G_i\gets 0\), where \(G_i=\partial_i\ellhat_T^{\rm FSGQ}\) is the
\(i\)-th component of the full score vector.
\item For \(t=1,\ldots,T\), symmetrize \(P_{t-1}\), then factor
\(P_{t-1}=C_{t-1}C_{t-1}^\top\) on the saved branch.  For this fixed
parameter component \(\theta_i\), compute
\(\dot C_{t-1}=DC(P_{t-1})[\dot P_{t-1}]\), and then run the prediction value
 equations \eqref{eq:p31-pred-placed-points}--\eqref{eq:p31-pred-cov} and
prediction sensitivity equations \eqref{eq:p31-dot-pred-place}--\eqref{eq:p31-dot-pred-cov}.
\item If the symmetrized predictive covariance \(P_t^-\) fails the declared
Cholesky or positive-definite test in \eqref{eq:p32-failure-eigen}, stop and
report a predictive-branch failure.
\item Factor \(P_t^-=C_t^-(C_t^-)^\top\) on the saved branch, compute
\(\dot C_t^-=DC(P_t^-)[\dot P_t^-]\), and then run the observation value
equations \eqref{eq:p31-obs-placed-points}--\eqref{eq:p31-fsgq-loglik}, which
produce the candidate one-step increment \(\ellhat_t^{\rm FSGQ}\), together with the
observation sensitivity equations \eqref{eq:p31-dot-obs-place}--\eqref{eq:p31-dot-Cxz}.
\item If the symmetrized innovation covariance \(S_t\) fails the declared
positive-definite branch test in \eqref{eq:p32-failure-eigen}, stop and report
an innovation-branch failure.
\item Solve \(S_tu_t=v_t\), so \(u_t\) is the innovation solve variable
representing the action of \(S_t^{-1}\) on \(v_t\) without forming an explicit
inverse.  The observation sensitivity stage has already
produced \(\dot{\bar z}_t\), hence \(\dot v_t=-\dot{\bar z}_t\), together with
\(\dot S_t\) on this same saved branch for this same parameter component
\(\theta_i\).  Compute the \(i\)-th one-step score component
\(\dot\ellhat_t^{\rm FSGQ}=\partial_i\ellhat_t^{\rm FSGQ}\) by
\eqref{eq:p31-score}, update the running derivative by
\(G_i\leftarrow G_i+\dot\ellhat_t^{\rm FSGQ}\), and after the innovation branch
has passed update the running scalar by
\(\ellhat_T^{\rm FSGQ}\leftarrow \ellhat_T^{\rm FSGQ}+\ellhat_t^{\rm FSGQ}\).
\item Update \(m_t,P_t,\dot m_t,\dot P_t\) using
\eqref{eq:p31-kalman-gain}--\eqref{eq:p31-filter-cov} and
\eqref{eq:p31-dot-K}--\eqref{eq:p31-dot-P} for this fixed component
\(\theta_i\).  Then apply the same symmetrize-then-veto rule to the updated
carried covariance: replace \(P_t\) by \((P_t+P_t^\top)/2\), and if the
symmetrized matrix is not on the declared positive-definite factor branch, stop
and report a carried-covariance failure.
\end{enumerate}
\end{minipage}}

\section{Implementation Contract}
\label{sec:p31-implementation-contract}

\begin{center}
\small
\begin{tabular}{p{0.18\linewidth}p{0.30\linewidth}p{0.16\linewidth}p{0.24\linewidth}}
\toprule
object & definition & shape & reuse \\
\midrule
\(\xi^{(r)}\) & standardized sparse-grid node & \(n_x\) & value and gradient \\
\(w_r\) & accumulated signed weight after duplicate merge & scalar & value and gradient \\
\(C_{t-1},C_t^-\) & saved covariance factors on declared branch & \(n_x\times n_x\) & value, gradient, diagnostics \\
\(\chi_{t-1}^{(r)}\) & previous placed point & \(n_x\) & transition value and derivative \\
\(a_t^{(r)}\) & transition image & \(n_x\) & prediction moments \\
\(\chi_t^{(r)}\) & predictive placed point & \(n_x\) & observation value and derivative \\
\(z_t^{(r)}\) & observation image & \(n_y\) & observation moments \\
\(\bar z_t,S_t,C_{xz,t}\) & Gaussian projection moments & \(n_y\), \(n_y\times n_y\), \(n_x\times n_y\) & likelihood and update \\
\(v_t,K_t,u_t\) & innovation, Gaussian projection gain, and innovation solve variable with \(S_tu_t=v_t\) & \(n_y\), \(n_x\times n_y\), \(n_y\) & likelihood, update, gradient \\
\(\dot \chi,\dot a,\dot z\) & point sensitivities & matching value objects & gradient only \\
\(\dot m,\dot P,\dot S,\dot C_{xz}\) & moment and filter sensitivities & matching value objects & gradient and next step \\
data record & \(y_{1:T}\), preprocessing, missing-data policy if any & metadata & value, gradient, finite difference \\
initial law & \(m_0(\theta),P_0(\theta),\dot m_0,\dot P_0\) & \(n_x,n_x\times n_x\) & value and gradient start \\
branch record & cloud, one-dimensional rules, merge tolerance, zero-weight rule, node ordering, preprocessing policy, initial-condition policy, declared Cholesky factor branch, symmetrize-then-veto rule, thresholds, accepted/failure stage-time pattern & metadata & same-scalar finite difference \\
\bottomrule
\end{tabular}
\end{center}

\subsection{Input, Output, Invariant, And Failure Contract}
\label{subsec:p32-io-contract}

The full report-level routine has inputs
\begin{equation}
\begin{aligned}
  \mathcal I_{\rm in}
  =
  (&y_{1:T},\theta,L,\{I_\ell\}_{\ell=1}^{L+b-1},
  m_0(\theta),P_0(\theta),\\
  &\dot m_0^{(1:p)}(\theta),\dot P_0^{(1:p)}(\theta),
  f_\theta,h_\theta,Q_\theta,R_\theta,\\
  &D_xf_\theta,\partial_{1:p}f_\theta,
  D_xh_\theta,\partial_{1:p}h_\theta,
  \dot Q_\theta^{(1:p)},\dot R_\theta^{(1:p)},\\
  &\text{merge tolerance},\text{factor map},\text{branch thresholds}).
\end{aligned}
  \label{eq:p32-input-contract}
\end{equation}
It returns either
\begin{equation}
  \left(
  \ellhat_T^{\rm FSGQ}(\theta),
  \nabla_\theta\ellhat_T^{\rm FSGQ}(\theta),
  \{m_t,P_t\}_{t=0}^T,
  \mathcal D_{\rm diag}
  \right)
  \label{eq:p32-output-success}
\end{equation}
or a failure record
\begin{equation}
  \mathcal F=(\hbox{time},\hbox{stage},\hbox{reason},\hbox{diagnostic values}).
  \label{eq:p32-output-failure}
\end{equation}
The invariants are:
\begin{align}
  \sum_{r=1}^M w_r&=1
  &&\hbox{up to construction tolerance},
  \label{eq:p32-invariant-weight-sum}\\
  P_t^-&=(P_t^-)^\top,\quad
  S_t=S_t^\top,\quad
  P_t=P_t^\top
  &&\hbox{after explicit symmetrization},
  \label{eq:p32-invariant-symmetry}\\
  S_t&\succ0
  &&\hbox{on every accepted likelihood step},
  \label{eq:p32-invariant-positive-S}\\
  \mathcal B_{\rm value}&=\mathcal B_{\rm gradient}
  &&\hbox{same saved branch}.
  \label{eq:p32-invariant-same-branch}
\end{align}
The routine fails rather than repairs if
\begin{equation}
  \lambda_{\min}(S_t)<\epsilon_S,\qquad
  \lambda_{\min}(P_t^-)<\epsilon_P,\qquad
  \lambda_{\min}(P_t)<\epsilon_P,
  \label{eq:p32-failure-eigen}
\end{equation}
or if the declared factor map cannot be evaluated on the symmetrized
covariance.  An implementation may add a smoothed repair, but then the repair
is a new scalar and needs a new derivative.

\subsection{One-Step Value-Path Contract}
\label{subsec:p32-one-step-value-contract}

At one accepted time step, the value path is the ordered map
\begin{equation}
\begin{aligned}
  (m_{t-1},P_{t-1},y_t,\theta,\mathcal C)
  &\longmapsto C_{t-1}
  \longmapsto \chi_{t-1}^{(r)}
  \longmapsto a_t^{(r)}
  \longmapsto (m_t^-,P_t^-) \\
  &\longmapsto C_t^-
  \longmapsto \chi_t^{(r)}
  \longmapsto z_t^{(r)}
  \longmapsto (\bar z_t,S_t,C_{xz,t}) \\
  &\longmapsto (v_t,\ellhat_t^{\rm FSGQ},K_t)
  \longmapsto (m_t,P_t).
\end{aligned}
  \label{eq:p32-one-step-value-map}
\end{equation}
The accepted map in \eqref{eq:p32-one-step-value-map} is conditioned on two
successful branch checks: the predictive covariance \(P_t^-\) must pass the
declared Cholesky and positive-definite test before observation-point
placement, and the innovation covariance \(S_t\) must pass the declared branch
test before the Gaussian projection update is accepted.  The accepted outputs
propagated to the next step are only
\((m_t,P_t)\) together with the global fixed branch record \(\mathcal B\).
Point clouds, moment arrays, gains, and diagnostic quantities may be cached for
gradient reuse, but they are not part of the carried filtering state itself.

\subsection{Default Numerical Constants}
\label{sec:p32-defaults}

The constants below are defaults for diagnostics, not claims of universal
validity:
\begin{center}
\small
\begin{tabular}{p{0.24\linewidth}p{0.22\linewidth}p{0.44\linewidth}}
\toprule
quantity & default & role \\
\midrule
duplicate merge tolerance & \(10^{-12}\max(1,\|\xi\|_\infty)\) & groups standardized nodes that should be identical under the declared rule \\
weight zero tolerance & \(10^{-14}\) & removes accumulated numerical-zero weights after merge \\
\(\epsilon_S\) & \(10^{-10}\) & veto threshold for innovation covariance positive definiteness \\
\(\epsilon_P\) & \(10^{-10}\) & veto threshold for carried and predictive covariance factors \\
\(\epsilon_\ell\) & \(10^{-3}\max(1,|\ellhat_T|)\) & pilot level-ladder change threshold \\
finite-difference relative tolerance & \(10^{-5}\) & same-scalar gradient diagnostic \\
maximum branch-invalid FD rows & 0 in final audit & branch changes mean the derivative is not being tested \\
\bottomrule
\end{tabular}
\end{center}
These values are deliberately conservative.  They are meant to reveal when the
fixed branch is fragile.  They are not tuning knobs for making a failing
branch look smooth.

\section{End-To-End Mathematical Algorithm}
\label{sec:p32-end-to-end-algorithm}

\fbox{\begin{minipage}{0.94\linewidth}
\textbf{FixedSGQF value and gradient, with all branch exits.}

\begin{enumerate}[leftmargin=0.22in]
\item Build the sparse-grid cloud from \(b=n_x\), level \(L\), and
\(\{I_\ell\}\).  Use \eqref{eq:p31-level-band},
\eqref{eq:p31-smolyak-coeff}, \eqref{eq:p32-smolyak-node},
\eqref{eq:p32-smolyak-tensor-weight}, and
\eqref{eq:p31-node-dictionary}.  Save
\(\mathcal C=\{(\xi^{(r)},w_r)\}_{r=1}^M\).
\item Check \eqref{eq:p32-invariant-weight-sum}.  If it fails, return
\(\mathcal F=(0,\hbox{cloud},\hbox{weight-sum failure},\cdots)\).
\item Initialize \(\ellhat_T^{\rm FSGQ}\gets0\),
\(G_i\gets0\) for \(i=1,\ldots,p\), and compute
\(m_0,P_0,\dot m_0^{(i)},\dot P_0^{(i)}\).
\item For \(t=1,\ldots,T\), symmetrize \(P_{t-1}\).  If the factor map
\(C(P_{t-1})\) fails, return a factor-branch failure.  Otherwise set
\(C_{t-1}=C(P_{t-1})\) and for each parameter component compute
\(\dot C_{t-1}^{(i)}=DC(P_{t-1})[\dot P_{t-1}^{(i)}]\).
\item Place previous points by \eqref{eq:p31-pred-placed-points}; compute
their derivatives by \eqref{eq:p31-dot-pred-place} componentwise in
\(i=1,\ldots,p\).
\item Evaluate transition values \eqref{eq:p31-transition-points} and
derivatives \eqref{eq:p31-dot-transition}.
\item Compute \(m_t^-\), \(P_t^-\), \(\dot m_t^{-(i)}\), and
\(\dot P_t^{-(i)}\) by \eqref{eq:p31-pred-mean},
\eqref{eq:p31-pred-cov}, \eqref{eq:p31-dot-pred-mean}, and
\eqref{eq:p31-dot-pred-cov}.  Symmetrize \(P_t^-\).  If
\eqref{eq:p32-failure-eigen} fails for \(P_t^-\), return a predictive
covariance failure.
\item Factor \(P_t^-\).  Set \(C_t^-=C(P_t^-)\) and compute
\(\dot C_t^{-(i)}=DC(P_t^-)[\dot P_t^{-(i)}]\) for each parameter component.
Place observation points by \eqref{eq:p31-obs-placed-points}; compute their
derivatives by \eqref{eq:p31-dot-obs-place}.
\item Evaluate observation values \eqref{eq:p31-obs-values} and derivatives
\eqref{eq:p31-dot-obs-value}.
\item Compute \(\bar z_t,\Delta_t^{(r)},S_t,C_{xz,t},v_t,\ellhat_t^{\rm FSGQ}\) by
\eqref{eq:p31-obs-mean}--\eqref{eq:p31-fsgq-loglik}.  Symmetrize \(S_t\).
If \(S_t\) fails \eqref{eq:p32-failure-eigen}, return an innovation
covariance failure.
\item Compute \(\dot{\bar z}_t^{(i)},\dot v_t^{(i)},\dot S_t^{(i)}\), and
\(\dot C_{xz,t}^{(i)}\) by \eqref{eq:p31-dot-zbar}--\eqref{eq:p31-dot-Cxz}.
\item Solve \(S_tu_t=v_t\).  Add \(\ellhat_t^{\rm FSGQ}\) to
\(\ellhat_T^{\rm FSGQ}\) only after the innovation branch has passed, and add
the componentwise score to \(G_i\) only on that same accepted innovation
branch.  Reuse this same solved \(u_t\) for every parameter component \(i\),
and add the componentwise score from \eqref{eq:p31-score} with
\(\dot v_t=\dot v_t^{(i)}\) and \(\dot S_t=\dot S_t^{(i)}\) to
\(G_i\).
\item Compute \(K_t\) first from \eqref{eq:p31-kalman-gain}.  Then for each
parameter component compute \(\dot K_t^{(i)}\) from \eqref{eq:p31-dot-K},
and update \(m_t,P_t,\dot m_t^{(i)},\dot P_t^{(i)}\) by
\eqref{eq:p31-filter-mean}--\eqref{eq:p31-filter-cov} and
\eqref{eq:p31-dot-m}--\eqref{eq:p31-dot-P}.  Symmetrize \(P_t\).  If the
carried covariance branch fails, return a carried covariance failure.
\item Return \((\ellhat_T^{\rm FSGQ},G_{1:p},\{m_t,P_t\},\mathcal D_{\rm diag})\).
\end{enumerate}
\end{minipage}}

The algorithm is deliberately mathematical rather than Pythonic.  Its purpose
is to say what a Python implementation must compute, save, and refuse to
change if it wants the reported gradient to be the derivative of the reported
value.

\subsection{Formula Inventory For The Value Path}
\label{subsec:p32-value-inventory}

For implementation review, the value path can be located in the report as
follows.
\begin{center}
\small
\begin{tabular}{p{0.34\linewidth}p{0.50\linewidth}}
\toprule
value-path item & report location \\
\midrule
symbol and dimension table & Section~\ref{subsec:p32-object-map} \\
exact filtering target & Section~\ref{sec:p31-ssm} \\
Gaussian projection moment equations & Section~\ref{sec:p31-gaussian-projection} \\
standardized sparse-grid cloud construction & Section~\ref{sec:p31-sgq-reconstruction} \\
merge convention and stored cloud semantics & Section~\ref{subsec:p32-cloud-construction} and Section~\ref{sec:p31-toy-grid} \\
per-step value inputs and outputs & Sections~\ref{sec:p31-value-path}, \ref{subsec:p32-io-contract}, and \ref{subsec:p32-one-step-value-contract} \\
prediction equations & \eqref{eq:p31-pred-placed-points}--\eqref{eq:p31-pred-cov} \\
observation and update equations & \eqref{eq:p31-obs-placed-points}--\eqref{eq:p31-filter-cov} \\
one-step log-likelihood increment & \eqref{eq:p31-fsgq-loglik} \\
implementation-order value/gradient algorithm & Section~\ref{sec:p32-end-to-end-algorithm} \\
branch checks and numerical conventions & Sections~\ref{subsec:p32-object-map}, \ref{subsec:p32-io-contract}, and \ref{sec:p32-defaults} \\
\bottomrule
\end{tabular}
\end{center}

\section{Finite-Difference Same-Scalar Check}
\label{sec:p31-fd}

The finite-difference check must call the same value routine with the same
saved branch \(\mathcal B\).  For parameter component \(i\), use
\begin{equation}
  D_i(h)
  =
  \frac{
  \ellhat_T^{\rm FSGQ}(\theta+he_i;\mathcal B)
  -
  \ellhat_T^{\rm FSGQ}(\theta-he_i;\mathcal B)}
  {2h}.
  \label{eq:p31-central-diff}
\end{equation}
Use the step ladder
\begin{equation}
  h\in\{10^{-1},10^{-2},10^{-3},10^{-4},10^{-5}\}\,
  \max\{1,|\theta_i|\}.
  \label{eq:p31-fd-ladder}
\end{equation}
For every \(h\), the \(+\!h\) and \(-\!h\) evaluations must reuse the same
saved structural branch record \(\mathfrak B\): the same cloud,
one-dimensional rules, preprocessing policy, initial-condition policy,
declared Cholesky factor branch, merge tolerance, zero-weight rule, node
ordering, and symmetrize-then-veto rule.  The numerical initial values
\(m_0(\theta\pm he_i),P_0(\theta\pm he_i)\) are recomputed by that same policy.
Each evaluation also returns its realized accepted/failure stage-time pattern.
A row is branch-valid only if both perturbations match the base evaluation in
that full branch-identity record and in the realized accepted/failure
stage-time pattern through every likelihood contribution included in the scalar
comparison.  Equal factor family but different first failure location counts as
branch mismatch.  Equal first failure location but different accepted-prefix
pattern also counts as branch mismatch.  If either side hits a different branch
or returns a veto, that row is marked branch-invalid and is not interpreted as a
derivative check.  For valid rows, compare \(D_i(h)\) against the analytical
score \(G_i\) by
\begin{equation}
  e_i(h)=
  \frac{|D_i(h)-G_i|}
  {\max\{1,|D_i(h)|,|G_i|\}}.
  \label{eq:p31-fd-relative-error}
\end{equation}
The diagnostic passes when at least two consecutive valid rows have
\(e_i(h)\le 10^{-5}\) and the smaller-step row is no worse than ten times the
larger-step row.  This is a numerical diagnostic, not a theorem: failure means
one of three things is likely.  The derivative formula is wrong, the value and
gradient paths are not using the same scalar, or the branch is too nonsmooth or
ill-conditioned at that point.

\paragraph{Implementation-facing FD pseudocode summary.}
A routine of the form \texttt{same\_scalar\_fd\_check(}$i,h_{\text{ladder}},\theta,\mathfrak B$\texttt{)}
should:
\begin{enumerate}[leftmargin=0.28in]
  \item evaluate the value path at \(\theta\pm he_i\) using the same saved
  structural branch metadata,
  \item mark a row branch-invalid if either side changes the structural branch
  or the realized accepted/failure stage-time pattern,
  \item compute the central difference only for branch-valid rows,
  \item compare that central difference against the analytical score with
  \eqref{eq:p31-fd-relative-error}, and
  \item report whether two consecutive valid rows satisfy the declared
  tolerance rule.
\end{enumerate}
This summary is the operational finite-difference contract for an
implementation; the equations above remain the authoritative mathematical
definition.

\paragraph{A scalar numerical trace.}
Let \(X\sim\N(0,1)\), \(h_\theta(x)=\theta x\), \(R=1\), and \(y=2\).  Any
fixed Gaussian rule with exact second moment gives
\begin{equation}
  \bar z=0,\qquad S(\theta)=1+\theta^2,\qquad
  \ellhat(\theta)
  =
  -\frac12\left\{\log(1+\theta^2)+\frac{4}{1+\theta^2}+\log(2\pi)\right\}.
  \label{eq:p31-fd-toy-likelihood}
\end{equation}
At \(\theta=1\),
\begin{equation}
  G
  =
  -\frac12
  \left\{
  \frac{2}{2}
  -
  \frac{4\cdot2}{2^2}
  \right\}
  =
  0.5.
  \label{eq:p31-fd-toy-gradient}
\end{equation}
Central differences give:
\begin{center}
\begin{tabular}{rrrr}
\toprule
\(h\) & \(D(h)\) & \(G\) & \(|D(h)-G|\) \\
\midrule
\(10^{-1}\) & 0.5008108250 & 0.5 & \(8.10825\times10^{-4}\) \\
\(10^{-2}\) & 0.5000083311 & 0.5 & \(8.33108\times10^{-6}\) \\
\(10^{-3}\) & 0.5000000833 & 0.5 & \(8.33331\times10^{-8}\) \\
\(10^{-4}\) & 0.5000000008 & 0.5 & \(8.33722\times10^{-10}\) \\
\bottomrule
\end{tabular}
\end{center}

\paragraph{A cloud-sensitive nonlinear trace.}
The previous trace tests the innovation-score algebra.  To exercise the fixed
cloud, use the two-dimensional level-2 cloud in Section~\ref{sec:p31-toy-grid}
and the nonlinear observation
\begin{equation}
  h_\theta(x_1,x_2)=\theta(x_1^2+x_2^2),
  \qquad
  X\sim\N(0,I_2),\qquad R=1,\qquad y=3.
  \label{eq:p31-cloud-sensitive-model}
\end{equation}
For the merged cloud in Section~\ref{sec:p31-toy-grid},
\begin{equation}
  \sum_r w_r(x_{1,r}^2+x_{2,r}^2)=2,
  \qquad
  \sum_r w_r(x_{1,r}^2+x_{2,r}^2)^2=6.
  \label{eq:p31-cloud-sensitive-moments}
\end{equation}
Therefore
\begin{equation}
  \bar z(\theta)=2\theta,\qquad
  S(\theta)=1+2\theta^2,\qquad
  v(\theta)=3-2\theta,
  \label{eq:p31-cloud-sensitive-likelihood-objects}
\end{equation}
if the parameter scales the whole nonlinear observation linearly.  To exercise
the parameter through placed nonlinear values more sharply, set instead
\begin{equation}
  h_\theta(x_1,x_2)=\theta^2(x_1^2+x_2^2),
  \label{eq:p31-cloud-sensitive-square-param}
\end{equation}
which gives
\begin{equation}
  \bar z(\theta)=2\theta^2,\qquad
  S(\theta)=1+2\theta^4,\qquad
  v(\theta)=3-2\theta^2.
  \label{eq:p31-cloud-sensitive-square-objects}
\end{equation}
At \(\theta=1\), the score from \eqref{eq:p31-score} is
\begin{equation}
  G
  =
  -\frac12
  \left[
    \frac{8}{3}
    +2(-4)\frac{1}{3}
    -\frac{1}{3^2}8
  \right]
  =
  \frac49.
  \label{eq:p31-cloud-sensitive-score}
\end{equation}
Central differences for the saved cloud scalar give:
\begin{center}
\begin{tabular}{rrrr}
\toprule
\(h\) & \(D(h)\) & \(G\) & \(|D(h)-G|\) \\
\midrule
\(10^{-1}\) & 0.5200303321 & 0.4444444444 & \(7.55859\times10^{-2}\) \\
\(10^{-2}\) & 0.4452146703 & 0.4444444444 & \(7.70226\times10^{-4}\) \\
\(10^{-3}\) & 0.4444521481 & 0.4444444444 & \(7.70369\times10^{-6}\) \\
\(10^{-4}\) & 0.4444445215 & 0.4444444444 & \(7.70368\times10^{-8}\) \\
\bottomrule
\end{tabular}
\end{center}
This example depends on the merged sparse-grid fourth moment in
\eqref{eq:p31-cloud-sensitive-moments}.  A duplicate-merge bug or a different
saved cloud changes \(S(\theta)\), \(G\), and the finite-difference table.

\section{Diagnostics, Accuracy, Memory, And Performance Tests}
\label{sec:p31-validation}

Fixed SGQF should be tested as an approximate deterministic likelihood, not as
an exact posterior method.  The minimum validation suite is:

\begin{enumerate}[leftmargin=0.24in]
\item \textbf{Construction checks.}  Verify duplicate-node accumulation,
weight totals, signed-weight counts, and deterministic cloud reproduction for
small \((b,L)\) cases such as Section~\ref{sec:p31-toy-grid}.
\item \textbf{Linear-Gaussian recovery.}  For affine \(f_\theta,h_\theta\),
the Gaussian projection recursion should match the Kalman filter up to
quadrature and floating-point tolerance.
\item \textbf{Quadratic scalar cell.}  Use \eqref{eq:p31-quadratic-cell} to
show both the moment accuracy and the posterior-shape limitation.
\item \textbf{Same-scalar finite differences.}  Run
\eqref{eq:p31-central-diff} for representative parameters and step sizes.
\item \textbf{Signed-weight stability.}  Record the minimum eigenvalues of
symmetrized \(P_t^-\), \(S_t\), and \(P_t\), plus any branch veto.
\item \textbf{Memory and runtime scaling.}  Record \(M_{b,L}\), model
evaluations, wall time, and peak memory as \(b,L,T\) change.
\item \textbf{Accuracy ladder.}  Compare fixed SGQF against tensor-product GHQ
in small dimension, Kalman truth in linear-Gaussian models, and particle or TT
diagnostics in nonlinear models where those references are available.
\end{enumerate}

Passing these checks supports the statement that the implementation evaluates
and differentiates the declared approximate scalar.  It does not prove exact
posterior accuracy.

\subsection{Mathematical Test Models}
\label{subsec:p32-test-models}

The validation suite should use models that expose different failure modes.
The following models are mathematical specifications for future code tests.  They
are ordered to mirror the report's argumentative burden: first recover exact or
near-exact reference behavior where possible, then expose the Gaussian-surrogate
limitation, then test the cloud-specific and high-dimensional claims that later
feed into the method-selection argument of Section~\ref{sec:p31-comparison}.

\paragraph{Model A: linear-Gaussian recovery.}
Let
\begin{align}
  X_t &= A_\theta X_{t-1}+\eta_t,
  &
  \eta_t&\sim\N(0,Q_\theta),
  \label{eq:p32-test-lg-state}\\
  Y_t &= H_\theta X_t+\varepsilon_t,
  &
  \varepsilon_t&\sim\N(0,R_\theta).
  \label{eq:p32-test-lg-obs}
\end{align}
The Kalman filter gives the exact likelihood and exact filtering moments.  A
correct FixedSGQF implementation should match the Kalman mean, covariance, and
innovation scalar up to floating-point tolerance because all required moments
are at most quadratic in the Gaussian coordinate.  This test checks moment
assembly, covariance factors, likelihood algebra, and derivative propagation,
but it does not stress nonlinear quadrature accuracy.

\paragraph{Model B: scalar quadratic observation.}
Use
\begin{align}
  X_t &= \rho X_{t-1}+\eta_t,
  &
  \eta_t&\sim\N(0,q),
  \label{eq:p32-test-quad-state}\\
  Y_t &= X_t^2+\varepsilon_t,
  &
  \varepsilon_t&\sim\N(0,r).
  \label{eq:p32-test-quad-obs}
\end{align}
This model is small enough to compare against dense numerical integration.
For small \(r\) and positive observations, the exact posterior can have two
lobes.  The test should report both moment agreement and posterior-shape
failure.  It is the guardrail against accidentally describing FixedSGQF as an
exact nonlinear posterior filter.

\paragraph{Model C: cloud-sensitive two-dimensional nonlinear observation.}
Let
\begin{align}
  X_t&=\rho X_{t-1}+\eta_t,\qquad
  \eta_t\sim\N(0,qI_2),
  \label{eq:p32-test-cloud-state}\\
  Y_t&=\theta^2(X_{t,1}^2+X_{t,2}^2)+\varepsilon_t,\qquad
  \varepsilon_t\sim\N(0,r).
  \label{eq:p32-test-cloud-obs}
\end{align}
The two-dimensional level-2 cloud in Section~\ref{sec:p31-toy-grid} gives the
explicit moments in \eqref{eq:p31-cloud-sensitive-moments}.  This test catches
duplicate-node and signed-weight errors because the fourth moment changes when
the cloud is assembled incorrectly.

\paragraph{Model D: block-local high-dimensional nonlinear model.}
For \(b=md\), partition \(x\in\R^b\) into \(m\) blocks \(x_{[j]}\in\R^d\).
Let
\begin{align}
  X_t&=A_\theta X_{t-1}+\eta_t,
  &
  \eta_t&\sim\N(0,Q_\theta),
  \label{eq:p32-test-block-state}\\
  Y_{t,j}
  &=
  \alpha_j^\top x_{t,[j]}
  +\beta_j\|x_{t,[j]}\|^2
  +\gamma_j x_{t,[j]}^\top B_jx_{t,[j+1]}
  +\varepsilon_{t,j},
  \label{eq:p32-test-block-obs}
\end{align}
with \(j=1,\ldots,m\) and \(x_{[m+1]}=x_{[1]}\).  This model has controlled
low-order interactions.  It is designed to test whether moderate sparse-grid
levels exploit interaction sparsity and whether point count, memory, and
runtime grow as expected.  Diagnostics should report \(b,L,M_{b,L},T\), model
evaluations per step, peak memory, wall time, and branch failures.

\paragraph{Model E: deliberately hard all-coordinate interaction.}
Let
\begin{equation}
  Y_t
  =
  \delta\prod_{j=1}^b(1+\alpha X_{t,j})
  +\varepsilon_t,
  \qquad \varepsilon_t\sim\N(0,r).
  \label{eq:p32-test-hard}
\end{equation}
This model is not expected to be easy for a moderate sparse grid.  It probes
the failure mode where high-order coordinate interactions matter.  A good
validation report should not hide this case; it should show when the level
ladder or covariance diagnostics reject the fixed grid.

\paragraph{Model F: Jia--Xin--Cheng orbit-style benchmark.}
Jia, Xin, and Cheng evaluate SGQF on an orbit estimation problem
\cite[Section 4]{jia2012sparsegrid}.  This note does not reproduce their full
experiment, but a future implementation should include an orbit-style
benchmark because it is the source paper's main empirical setting.  The test
should record the state dimension, observation geometry, noise covariances,
chosen sparse-grid level, point count, RMSE, likelihood diagnostics, and
finite-difference branch status.  This benchmark supports source-faithful
comparison; it is not a theorem of general performance.

\subsection{Large-Scale Report Template}
\label{subsec:p32-large-scale-report}

For each large-scale run, record
\begin{equation}
\begin{aligned}
  \mathcal R
  =
  (&b,L,M_{b,L},T,\hbox{number of parameters},\hbox{random seed},\\
  &\hbox{model evaluations per step},\hbox{peak memory},\hbox{wall time},\\
  &\min_t\lambda_{\min}(P_t^-),
  \min_t\lambda_{\min}(S_t),
  \min_t\lambda_{\min}(P_t),\\
  &\ellhat_T,\|\nabla\ellhat_T\|,
  \hbox{finite-difference table},\hbox{branch failures}).
\end{aligned}
  \label{eq:p32-large-scale-report}
\end{equation}
Accuracy claims should be made only against a named comparator:
\begin{equation}
  \hbox{error}
  =
  \begin{cases}
  |\ellhat_T-\ell_T^{\rm Kalman}|, & \hbox{linear-Gaussian test},\\
  |\ellhat_T-\ell_T^{\rm dense}|, & \hbox{small dense quadrature test},\\
  \hbox{RMSE or posterior diagnostic against PF/TT reference}, &
  \hbox{nonlinear large test}.
  \end{cases}
  \label{eq:p32-comparator-errors}
\end{equation}
Memory and performance do not rescue a failed accuracy or branch diagnostic.
They are interpreted after the veto checks.  These validation objects therefore
set the stage for the final two interpretive tasks of the report: first, to say
what adaptive sparse-grid selection can still contribute once same-scalar
requirements are imposed, and second, to decide how FixedSGQF should be judged
against neighboring algorithmic lanes.

\section{Adaptive Sparse Grids As Grid Design}
\label{sec:p31-adaptive}

Singh et al. \cite{singh2018adaptive} choose sparse-grid indices adaptively by
using admissible index sets, local error indicators, active and old index
sets, a global error estimate, and a tolerance.  That is useful for designing a
cloud before inference.  For HMC, however, the live target cannot silently
change its grid as \(\theta\) moves.  The safe interpretation is
\begin{equation}
  \text{adaptive pilot chooses a cloud}
  \quad\Longrightarrow\quad
  \text{freeze the cloud}
  \quad\Longrightarrow\quad
  \text{run FixedSGQF}.
  \label{eq:p31-adaptive-to-fixed}
\end{equation}
If the adaptive index set is rerun inside the likelihood, the value path is a
piecewise-defined algorithm with discrete changes.  The formulas in
Section~\ref{sec:p31-gradient} do not differentiate those changes.  This is why
adaptive sparse-grid selection enters the note as an upstream design aid rather
than as the selected differentiable likelihood lane.  With that restriction now
explicit, the report can finally compare FixedSGQF against neighboring methods on
the exact criteria that matter for selection.

\section{Relation To Neighboring High-Dimensional Filtering Proposals}
\label{sec:p31-comparison}

The previous sections have fixed the mathematical object this note is trying to
carry: one Gaussian surrogate \((m_t,P_t)\) together with the deterministic
same-scalar likelihood
\begin{equation}
  \ellhat_T^{\rm FSGQ}(\theta)
  =
  \sum_{t=1}^T \ellhat_t^{\rm FSGQ}(\theta).
  \label{eq:p32-selection-fixedsgqf-scalar}
\end{equation}
This section is therefore not a general survey.  It asks a narrower mathematical
question: among nearby approximate filtering lanes, which constructions define a
usable deterministic surrogate scalar, which preserve an analytical derivative
of that same scalar on a fixed branch, and which avoid the dense-growth burden
that would make the lane unusable in the intended high-dimensional regime?

\paragraph{Quick comparison table.}
\begin{center}
\small
\begin{tabular}{p{0.16\linewidth}p{0.19\linewidth}p{0.19\linewidth}p{0.17\linewidth}p{0.20\linewidth}}
\toprule
method & carried object & moment engine & same-scalar derivative lane & main limitation \\
\midrule
EKF & one Gaussian & local linearization & yes, for the linearized scalar & weak nonlinear moment fidelity \\
UKF/CKF & one Gaussian & fixed low-order deterministic points & yes & lower-order deterministic rule \\
Dense GHQF & one Gaussian & dense tensor-product Gaussian quadrature & yes & exponential point growth \\
FixedSGQF & one Gaussian & fixed sparse-grid quadrature & yes & still one Gaussian; signed-weight branch sensitivity \\
Adaptive sparse grid & one Gaussian & parameter-dependent active grid & only if frozen first & target changes when the live grid changes \\
Particle / TT lanes & richer non-Gaussian object & samples or density representation & different target class & heavier approximation object, different contract \\
\bottomrule
\end{tabular}
\end{center}

For any deterministic Gaussian-surrogate lane \(M\), write its induced scalar as
\begin{equation}
  \ell_T^{(M)}(\theta)
  =
  \sum_{t=1}^T
  \log \mathcal N\bigl(y_t;\bar z_t^{(M)}(\theta),S_t^{(M)}(\theta)\bigr).
  \label{eq:p32-selection-generic-scalar}
\end{equation}
and its moment engine as the operator
\begin{equation}
  (\bar z_t^{(M)},S_t^{(M)},C_{xz,t}^{(M)})
  =
  \mathcal Q_M\bigl(f_\theta,h_\theta,m_{t-1},P_{t-1},Q_\theta,R_\theta\bigr).
  \label{eq:p32-selection-generic-operator}
\end{equation}
The comparison is then mathematical rather than rhetorical: nearby Gaussian
filters differ by the operator \(\mathcal Q_M\) that feeds the same Gaussian
update algebra, while particle and tensor-density methods differ already at the
level of the carried object.

For a deterministic point rule with nodes \(\xi_r^{(M)}\) and weights
\(w_r^{(M)}\), that operator has the explicit form
\begin{equation}
  \bar z_t^{(M)}
  \approx
  \sum_{r=1}^{M_M} w_r^{(M)}
  \,h_\theta\!\bigl(m_t^-+C_t^-\xi_r^{(M)}\bigr),
  \label{eq:p32-selection-zbar}
\end{equation}
\begin{equation}
  S_t^{(M)}
  \approx
  R_\theta+
  \sum_{r=1}^{M_M} w_r^{(M)}
  \Delta_{t,r}^{(M)}\Delta_{t,r}^{(M)\top},
  \qquad
  \Delta_{t,r}^{(M)}:=h_\theta(m_t^-+C_t^-\xi_r^{(M)})-\bar z_t^{(M)},
  \label{eq:p32-selection-S}
\end{equation}
\begin{equation}
  C_{xz,t}^{(M)}
  \approx
  C_t^-\!\left(
  \sum_{r=1}^{M_M} w_r^{(M)}\,
  \xi_r^{(M)}\Delta_{t,r}^{(M)\top}
  \right).
  \label{eq:p32-selection-Cxz}
\end{equation}
Thus the decisive question is not whether these methods are all “Gaussian,” but
what kind of moment engine they use and what that engine implies for the scalar
and its derivative.

\paragraph{What disqualifies a lane for this note?}
A lane is not selected here if any one of the following fails:
\begin{equation}
\begin{aligned}
  &\text{(i) a declared deterministic approximate scalar objective,}\\
  &\text{(ii) an analytical derivative of that same scalar on a fixed smooth branch,}\\
  &\text{(iii) a representation that avoids dense tensor-product growth in the intended regime.}
\end{aligned}
  \label{eq:p32-selection-three-conditions}
\end{equation}
This is not a theorem of universal method superiority.  It is the local choice
architecture of this report.  Conditions (i)--(iii) are used here as hard vetoes:
a method that fails one of them is not selected for this lane.  Once those
conditions are met, the later comparisons concern which eligible lane best
realizes the report's strong preferences.

\paragraph{EKF and local linearization.}
The extended Kalman filter carries the same kind of Gaussian object
\((m_t,P_t)\), but its moment engine is not
\eqref{eq:p32-selection-zbar}--\eqref{eq:p32-selection-Cxz}.  Instead it
substitutes local linear surrogates,
\begin{equation}
  f_\theta(x)\approx f_\theta(m)+F_\theta(x-m),
  \qquad
  h_\theta(x)\approx h_\theta(m^-)+H_\theta(x-m^-),
  \label{eq:p32-selection-ekf-linearization}
\end{equation}
so that the predictive and observation moments are inherited from Jacobian
closure rather than from explicit nonlinear point evaluation.  In particular,
the moment operator becomes
\begin{equation}
  \mathcal Q_{\mathrm{EKF}}
  \approx
  \bigl(h_\theta(m_t^-),\ H_\theta P_t^- H_\theta^\top +R_\theta,\ P_t^-H_\theta^\top\bigr),
  \label{eq:p32-selection-ekf-operator}
\end{equation}
which is computationally cheap but sacrifices the explicit nonlinear pointwise
moment information preserved by \eqref{eq:p32-selection-zbar}--\eqref{eq:p32-selection-Cxz}.  The present
report does not reject EKF as incoherent; it rejects it because the selected lane
insists on evaluating the nonlinear maps themselves under a declared cloud.

\paragraph{Low-order sigma-point and cubature lanes.}
UKF and CKF remain viable deterministic Gaussian alternatives because they also
construct a scalar of the form \eqref{eq:p32-selection-generic-scalar}.  Their
difference lies in the point design,
\begin{equation}
  \mathcal Q_{\mathrm{low}}
  \sim
  \sum_{r=1}^{M_{\mathrm{low}}} \tilde w_r F(\tilde \xi_r),
  \qquad
  M_{\mathrm{low}}=O(n_x),
  \label{eq:p32-selection-low-order}
\end{equation}
where the low point count is attractive but the deterministic rule remains a
lower-order moment design.  FixedSGQF instead chooses the sparse-grid rule
\begin{equation}
  \mathcal Q_{\mathrm{SGQF}}
  \sim
  \sum_{r=1}^{M_{b,L}} w_r F(\xi^{(r)}),
  \qquad
  M_{b,L}=O\!\bigl(b^{L-1}\bigr)
  \quad\text{for fixed }L,
  \label{eq:p32-selection-sparse-rule}
\end{equation}
The mathematical tradeoff is therefore visible in two dimensions at once:
\begin{equation}
  M_{\mathrm{low}}=O(n_x)
  \qquad\text{versus}\qquad
  M_{b,L}=O\!\bigl(b^{L-1}\bigr),
  \label{eq:p32-selection-pointcount-contrast}
\end{equation}
while the corresponding moment operators satisfy
\begin{equation}
  \mathcal Q_{\mathrm{low}}
  \neq
  \mathcal Q_{\mathrm{SGQF}}
  \quad\text{except in special low-order cases where the same moments are matched.}
  \label{eq:p32-selection-moment-contrast}
\end{equation}
So the comparison is not “deterministic versus nondeterministic,” but low-order
low-point deterministic moment closure versus higher-structure sparse-grid
deterministic moment closure.  The selected lane is willing to pay more than
\(O(n_x)\) points in exchange for a richer deterministic moment structure while
still avoiding the tensor-product explosion of dense GHQF.  The decisive point
for this note is not that low-order sigma-point rules are invalid, but that once
the project asks for a deterministic Gaussian-surrogate lane with a visibly
recorded cloud, merge, and branch identity, the sparse-grid construction carries
more of the moment-engine burden in the object itself rather than in a smaller
rule whose adequacy must be argued indirectly.

\paragraph{Standard and tensor-product GHQF.}
Dense Gaussian quadrature remains the cleanest low-dimensional reference, but it
fails the present lane in the intended regime because the point count grows as
\begin{equation}
  M_{\mathrm{tensor}}(b,L)=s_L^{\,b},
  \label{eq:p32-selection-tensor-growth}
\end{equation}
where \(s_L\) is the one-dimensional point count at the chosen level.  By
contrast, the sparse-grid cloud in this report has
\begin{equation}
  M_{\mathrm{SGQF}}(b,L)=M_{b,L}=O\!\bigl(b^{L-1}\bigr)
  \qquad\text{for fixed }L.
  \label{eq:p32-selection-sparse-growth}
\end{equation}
The note does not claim that sparse-grid quadrature is uniformly more accurate
than dense GHQF.  It claims something narrower: if the lane is defined by
deterministic Gaussian-surrogate likelihood evaluation in a high-dimensional
regime, then dense GHQF is too expensive to be the selected lane and should
remain a comparator or small-dimension reference.

\paragraph{Live adaptive sparse-grid selection.}
Adaptive sparse-grid methods are still useful, but not as the selected
same-scalar derivative lane.  If the active index set changes with \(\theta\),
then the scalar itself becomes
\begin{equation}
  \ell_T^{\mathrm{adap}}(\theta)
  =
  \sum_{t=1}^T \ell_t\bigl(\theta,\mathcal I_t(\theta)\bigr),
  \label{eq:p32-selection-adaptive-scalar}
\end{equation}
where \(\mathcal I_t(\theta)\) is the parameter-dependent active index set.  Its
formal derivative is then only local to regions where the index set is constant:
\begin{equation}
  \partial_i \ell_T^{\mathrm{adap}}(\theta)
  =
  \sum_{t=1}^T
  \partial_i \ell_t\bigl(\theta,\mathcal I_t(\theta)\bigr)
  \quad\text{only on regions where }\mathcal I_t(\theta)\text{ is locally constant.}
  \label{eq:p32-selection-adaptive-derivative}
\end{equation}
That is a different target class from the fixed scalar of this note.  The role of
adaptation here is therefore upstream cloud design, not the selected downstream
same-scalar likelihood lane.

\paragraph{Particle and tensor-density lanes.}
Particle and TT methods differ already at the level of the carried object.
Particle filtering carries weighted samples; squared TT carries a richer
non-Gaussian density approximation.  In schematic form,
\begin{equation}
  \text{particle lane: carried object } \{x_t^{(r)},\omega_t^{(r)}\}_{r=1}^N,
  \qquad
  \text{TT lane: carried object } \widehat p_t^{\mathrm{TT}}(x),
  \label{eq:p32-selection-richer-objects}
\end{equation}
whereas FixedSGQF carries only \((m_t,P_t)\) together with the surrogate scalar.
These methods are not ruled out because they are weak.  They are retained as
complementary lanes because they preserve posterior structure that one Gaussian
surrogate cannot preserve.  The present note is therefore not selecting the
method with the broadest posterior-shape fidelity; it is selecting the
deterministic Gaussian-surrogate lane with the cleanest same-scalar
likelihood-and-gradient contract.

\paragraph{Mathematical selection summary.}
The comparison therefore closes not with a slogan but with the following local
selection statement:
\begin{equation}
  \text{Selected lane}
  =
  \text{deterministic Gaussian-surrogate scalar}
  +
  \text{fixed-branch analytical derivative}
  +
  \text{sparse-grid moment engine with }M_{b,L}\ll M_{\mathrm{tensor}}.
  \label{eq:p32-selection-summary}
\end{equation}
This equation is not a theorem of universal optimality.  It is the summary of
the local choice architecture built in this report.  FixedSGQF is selected
because, among nearby deterministic Gaussian candidates, it is the lane that
jointly preserves: a declared deterministic scalar, an analytical derivative of
that same scalar on a fixed branch, an explicit approximation boundary, and a
cost story less explosive than dense GHQF while more stable than live
adaptive-grid selection.  That is the precise sense in which this note chooses
FixedSGQF.

\section{Where Each Construction Comes From}
\label{sec:p31-source-map}

\begin{center}
\small
\begin{tabular}{p{0.34\linewidth}p{0.28\linewidth}p{0.28\linewidth}}
\toprule
construction & support & use in this note \\
\midrule
state-space and conceptual filtering equations & \citet[Section 2]{jia2012sparsegrid} & Sections~\ref{sec:p31-ssm}--\ref{sec:p31-gaussian-projection} \\
Gaussian approximation filter moment structure & \citet[Eq. 15--23]{jia2012sparsegrid} plus Proposition~\ref{prop:p31-affine-projection} & moment projection and update equations \\
one-dimensional and tensor-product Gaussian quadrature & \citet[Section 2.2.1]{jia2012sparsegrid} & Section~\ref{sec:p31-ghq} \\
sparse-grid formula and index band & \citet[Eq. 26--29]{jia2012sparsegrid} & Section~\ref{sec:p31-sgq-reconstruction} \\
point/weight generation and duplicate accumulation & \citet[Algorithm 1]{jia2012sparsegrid} & fixed cloud \(\mathcal C\) and toy grid \\
quadrature exactness and UKF relation & \citet[Theorems 3.1--3.2]{jia2012sparsegrid} & source-scope orientation, not posterior guarantee \\
adaptive sparse-grid mechanics & \citet[Section 3]{singh2018adaptive} & offline grid-design interpretation \\
fixed scalar and analytical gradient & BayesFilter project derivation & Sections~\ref{sec:p31-same-scalar}--\ref{sec:p31-gradient} \\
comparison with squared TT & Zhao--Cui companion note and \citet{zhao2024ttsequential} & Section~\ref{sec:p31-comparison} \\
\bottomrule
\end{tabular}
\end{center}

\section{Notation And Formula Inventory}
\label{sec:p32-final-inventory}

For implementation review, the note can be navigated by the following inventory.
\begin{center}
\small
\begin{tabular}{p{0.33\linewidth}p{0.53\linewidth}}
\toprule
item & report location \\
\midrule
scientific problem and lane choice & Sections~\ref{sec:p32-problem-lane}--\ref{sec:p32-one-page} \\
exact object versus approximate object & \eqref{eq:p32-exact-vs-approximate-summary}, Section~\ref{sec:p31-computes} \\
global notation and dimensions & Section~\ref{subsec:p32-object-map} \\
Gaussian projection equations & Section~\ref{sec:p31-gaussian-projection} \\
source-order sparse-grid construction & Section~\ref{sec:p31-sgq-reconstruction} \\
merge convention and stored cloud semantics & Section~\ref{subsec:p32-cloud-construction}, Section~\ref{sec:p31-toy-grid} \\
value-path recursion & Section~\ref{sec:p31-value-path}, Section~\ref{subsec:p32-one-step-value-contract}, Section~\ref{subsec:p32-value-inventory} \\
fixed scalar and saved-branch contract & Section~\ref{sec:p31-same-scalar}, \eqref{eq:p31-fixed-scalar}, \eqref{eq:p32-branch-identity-record} \\
factor derivative convention & Section~\ref{sec:p31-gradient}, subsection `Square-Root Branch', equations \eqref{eq:p31-chol-A}--\eqref{eq:p31-chol-derivative} \\
gradient dependency chain and roles & \eqref{eq:p32-gradient-chain}, \eqref{eq:p32-score-role}, \eqref{eq:p32-propagation-role} \\
prediction and observation sensitivities & \eqref{eq:p31-dot-pred-place}--\eqref{eq:p31-dot-Cxz} \\
innovation-score derivative & Proposition~\ref{prop:p31-innovation-score}, \eqref{eq:p31-score} \\
posterior sensitivity propagation & \eqref{eq:p32-posterior-sensitivity-order}, \eqref{eq:p31-dot-K}--\eqref{eq:p31-dot-P} \\
one-component summary algorithm & Section~\ref{sec:p31-boxed-algorithm} \\
full implementation-order algorithm & Section~\ref{sec:p32-end-to-end-algorithm} \\
worked numeric oracle & Section~\ref{sec:p32-worked-example} \\
finite-difference same-scalar diagnostics & Section~\ref{sec:p31-fd} \\
neighboring-method positioning & Section~\ref{sec:p31-comparison} \\
validation models and report template & Section~\ref{sec:p31-validation}, Section~\ref{subsec:p32-large-scale-report} \\
\bottomrule
\end{tabular}
\end{center}

\section{Conclusion}
\label{sec:p31-conclusion}

This report has addressed a specific scientific problem: high-dimensional
nonlinear filtering in a regime where exact posterior propagation is not a
practical computational object, but where a deterministic and analytically
differentiable surrogate may still be scientifically valuable.  The proposal
advanced here is FixedSGQF, a lane that carries one Gaussian surrogate,
approximates its required nonlinear moments by a fixed sparse-grid cloud, and
differentiates the resulting deterministic same-scalar likelihood on a declared
branch.

What the report has constructed is therefore precise.  It has reconstructed the
Gaussian-projection sparse-grid filtering value path in source order, made the
saved branch explicit, given a branch-conditioned analytical gradient of the
same scalar that the value routine evaluates, provided a compact worked numeric
oracle, and written the value and gradient recursions in implementation order.
A reader should now be able to identify the exact carried state, the exact
accumulated scalar, the exact derivative contract, and the exact places where a
branch veto occurs.

The report has also been explicit about what it does \emph{not} claim.  It does
not claim exact nonlinear likelihood evaluation.  It does not claim that
sparse-grid exactness for selected Gaussian moments implies exact non-Gaussian
posterior shape.  It does not claim that one carried Gaussian can preserve
multimodality, heavy asymmetry, or all-coordinate interaction structure whenever
those features matter scientifically.  It does not replace richer non-Gaussian
density-approximation lanes such as Zhao--Cui squared TT.

The strongest scientific case for FixedSGQF is therefore not generic novelty but
a specific combination of properties.  The cloud is declared, the
covariance-factor convention is declared, the branch-veto logic is declared, the
scalar is deterministic on that branch, and the analytical gradient
differentiates that same declared scalar.  Those properties are exactly the
criteria the report used to distinguish FixedSGQF from its nearest deterministic
Gaussian alternatives: more explicit nonlinear point evaluation than EKF, a more
transparent sparse-grid cloud contract than UKF/CKF, a less explosive cost story
than dense GHQF, and a cleaner same-scalar derivative contract than live adaptive
grid selection.  This makes the method well suited to settings where
reproducibility, inspectability, and gradient-based parameter learning matter,
and where the filtering task can tolerate a Gaussian-surrogate carried object.
In regimes with lower effective interaction order, the sparse-grid moment engine
can also be scientifically plausible without requiring a full tensor-product
explosion.

The main caution is equally clear.  FixedSGQF should not be oversold as a full
non-Gaussian posterior method.  When posterior shape itself is central --- for
example under strong multimodality, hard global interactions, or branch
fragility --- the deterministic Gaussian-surrogate lane becomes scientifically
too narrow, and richer sample-based or density-based methods deserve priority.
That is not a defect in the present report; it is the explicit boundary that
keeps the proposal honest.

A fair review-panel judgment should therefore be the following.  FixedSGQF is a
coherent, mathematically explicit, self-contained, and implementation-ready
high-dimensional filtering lane.  It is narrower than a full non-Gaussian
density approximation, but that narrowness buys a transparent deterministic
surrogate and a branch-conditioned analytical gradient that can be audited,
implemented, and checked directly.  The approval claim remains bounded: this
report is not asserting that FixedSGQF is the best nonlinear filter overall, but
that it is the selected deterministic Gaussian-surrogate lane for the explicit
approximate likelihood-and-gradient objective developed here.  On that basis,
this report is sufficiently thorough and persuasive to justify approval of
FixedSGQF as one serious lane in a broader high-dimensional filtering program.

\bibliographystyle{plainnat}
\bibliography{../references}

\end{document}
